\chapter{Classical Decision Theory}

\section{Lecture 1, 19: Basic Concentration Inequalities}

We start with the introduction to concentration inequalities.

There are several refinements to the Chebyshev inequality. For many random variables, the moment generating function will exist in a neighborhood around 0. In these cases, we can use the mgf to produce a tail bound.

Define $\mu=\mathbb{E}[X]$. For any $t>0$, we have that,
\begin{align*}
\mathbb{P}((X-\mu) \geq u)=\mathbb{P}(\exp (t(X-\mu)) \geq \exp (t u)) \leq \frac{\mathbb{E}[\exp (t(X-\mu))]}{\exp (t u)}
\end{align*}
by applying Markov's inequality. Now $t$ is a parameter we can choose to get a tight upper bound, i.e., we can write this bound as:
\begin{align*}
\mathbb{P}((X-\mu) \geq u) \leq \inf _{0 \leq t \leq b} \exp (-t(u+\mu)) \mathbb{E}[\exp (t X)]
\end{align*}
This bound is known as Chernoff's bound. Suppose that, $X \sim N\left(\mu, \sigma^2\right)$, then a simple calculation gives that the mgf of $X$ is:
\begin{align*}
M_X(t)=\mathbb{E}[\exp (t X)]=\exp \left(t \mu+t^2 \sigma^2 / 2\right) .
\end{align*}
The mgf is defined for all $t$. To apply the Chernoff bound we then need to compute:
\begin{align*}
\inf _{t \geq 0} \exp (-t(u+\mu)) \exp \left(t \mu+t^2 \sigma^2 / 2\right)=\inf _{t \geq 0} \exp \left(-t u+t^2 \sigma^2 / 2\right),
\end{align*}
which is minimized when $t=u / \sigma^2$ which in turn yields the tail bound,
\begin{align*}
\mathbb{P}(X-\mu \geq u) \leq \exp \left(-u^2 /\left(2 \sigma^2\right)\right)
\end{align*}
This is often referred to as a one-sided or upper tail bound. We also have the two-sided Gaussian tail bound:
\begin{align*}
\mathbb{P}(|X-\mu| \geq u) \leq 2 \exp \left(-u^2 /\left(2 \sigma^2\right)\right)
\end{align*}

The main thing to observe is that this inequality is much sharper than Chebyshev's inequality. In particular, suppose we consider the average of i.i.d. Gaussian random variables, i.e. we have $X_1, \ldots, X_n \sim N\left(\mu, \sigma^2\right)$ and we construct the estimate:
\begin{align*}
\widehat{\mu}=\frac{1}{n} \sum_{i=1}^n X_i
\end{align*}
Using the fact that the average of Gaussian RVs is Gaussian, we obtain that $\widehat{\mu}$ has a $N\left(\mu, \sigma^2 / n\right)$ distribution. In this case, using the Gaussian tail bound we derived, we obtain that,
\begin{align*}
\mathbb{P}(|\widehat{\mu}-\mu| \geq t \sigma / \sqrt{n}) \leq 2 \exp \left(-t^2 / 2\right)
\end{align*}
This is an example of an exponential tail inequality. It turns out that the Gaussian tail inequality from the previous section is much more broadly applicable to a class of random variables called sub-Gaussian random variables. 

\begin{de}
    A random variable $X$ with mean $\mu$ is sub-Gaussian if there exists a positive number $\sigma$ such that,
\begin{align*}
\mathbb{E}[\exp (t(X-\mu))] \leq \exp \left(\sigma^2 t^2 / 2\right),
\end{align*}
for all $t \in \mathbb{R}$. Such a random variable is called $\sigma$-sub-Gaussian.
\end{de}
\begin{remark}
    Gaussian random variables with variance $\sigma^2$ satisfy the above condition with equality, so a $\sigma$-sub-Gaussian random variable basically just has an mgf that is dominated by a Gaussian with variance $\sigma$.
\end{remark}

\begin{eg}
    (Rademacher random variables are 1-sub-Gaussian) Let $X$ be a Rademacher random variable, which takes the values $\{+1,-1\}$ equiprobably. In this case, we can see that,
\begin{align*}
\begin{aligned}
\mathbb{E}[\exp (t X)] & =\frac{1}{2}[\exp (t)+\exp (-t)] \\
& =\frac{1}{2}\left[\sum_{k=0}^{\infty} \frac{(-t)^k}{k!}+\sum_{k=0}^{\infty} \frac{t^k}{k!}\right] \\
& =\sum_{k=0}^{\infty} \frac{t^{2 k}}{(2 k)!} \leq \sum_{k=0}^{\infty} \frac{t^{2 k}}{2^k k!} \\
& =\exp \left(t^2 / 2\right)
\end{aligned}
\end{align*}
This shows that Rademacher random variables are 1-sub-Gaussian.
\end{eg}

\begin{eg}
    (Bounded random variables are sub-Gaussian) Let $X$ be a random variable with zero mean and with support on some bounded interval $[a, b]$. WLOG, one may assume that $X$ has a zero mean. Let $X^{\prime}$ denote an independent copy of $X$ then we have that,
\begin{align*}
\mathbb{E}_X[\exp (t X)]=\mathbb{E}_X\left[\exp \left(t\left(X-\mathbb{E}\left[X^{\prime}\right]\right)\right] \leq \mathbb{E}_{X, X^{\prime}}\left[\exp \left(t\left(X-X^{\prime}\right)\right],\right.\right.
\end{align*}
using Jensen's inequality and the convexity of the function $g(x)=\exp (x)$.
Now, let $\epsilon$ be a Rademacher random variable. Then note that the distribution of $X-X^{\prime}$ is identical to the distribution of $X^{\prime}-X$ and, more importantly, of $\epsilon\left(X-X^{\prime}\right)$. So we obtain that,
\begin{align*}
\begin{aligned}
\mathbb{E}_{X, X^{\prime}}\left[\exp \left(t\left(X-X^{\prime}\right)\right]\right. & =\mathbb{E}_{X, X^{\prime}}\left[\mathbb{E}_\epsilon\left[\exp \left(t \epsilon\left(X-X^{\prime}\right)\right]\right]\right. \\
& \leq \mathbb{E}_{X, X^{\prime}}\left[\exp \left(t^2\left(X-X^{\prime}\right)^2 / 2\right]\right.
\end{aligned}
\end{align*}

where we just use the result that Rademacher random variables are 1-sub-Gaussian. Now using the boundedness of $X-X'$, we obtain that
\begin{align*}
\mathbb{E}_X[\exp (t X)] \leq \exp \left(t^2(b-a)^2 / 2\right),
\end{align*}
which in turn shows that bounded random variables are $(b-a)$-sub-Gaussian. Using some refined analysis, one can show that $X$ is indeed \textcolor{blue}{$(b-a)/2$-sub-Gaussian.}
\end{eg}
 

It is straightforward to go through the above Chernoff bound to conclude that for a sub-Gaussian random variable, we have the same two-sided exponential tail bound.

\begin{thm}
Let $X$ be a $\sigma$-sub-Gaussian random variable, then
\begin{align*}
\mathbb{P}(|X-\mu| \geq u) \leq 2 \exp \left(-u^2 /\left(2 \sigma^2\right)\right)
\end{align*}
\end{thm}



The following elementary result will be useful.
\begin{lem}
    Suppose that $X=\left(X_1, \ldots, X_n\right)$ is a random vector such that $X_i$ is $\sigma_i$-sub-Gaussian for $i=1, \ldots, n$. If $X_i$ are independent and $u \in \mathbb{R}^n$, then the random variable $u^{\top} X$ is sub-Gaussian with parameter $\sigma=\sqrt{\sigma_1^2 u_1^2+\ldots+\sigma_n^2 u_n^2}$.
\end{lem}

We list a couple of useful observations that follow directly from this lemma. They show that the $\sigma$ parameter behaves similarly to the standard deviation.

\begin{cor}
    For a collection of independent $\sigma$-sub-Gaussian random variables $X_1, \ldots, X_n$, their average $\bar{X}_n$ is sub-Gaussian with parameter $\sigma / \sqrt{n}$.
\end{cor}

\begin{cor}
    If $X$ is $\sigma$-sub-Gaussian, then $a X+b(a \in \mathbb{R}, b \in \mathbb{R})$ is sub-Gaussian with parameter $|a| \sigma$. In particular, $-X$ is $\sigma$-sub-Gaussian.
\end{cor}

\begin{cor}
    If $X=\left(X_1, \ldots, X_n\right)$ is a vector of independent random variables such that each $X_i$ is $\sigma$-sub-Gaussian, and $\|\boldsymbol{u}\|=1$, then $\boldsymbol{u}^T X$ is $\sigma$-sub-Gaussian.
\end{cor}



\begin{cor}
    Suppose we have $n$ i.i.d. $\sigma$-sub-Gaussian random variables $X_1, X_2, \ldots, X_n$ with a common mean $\mu$. Then we also have the following tail estimate for their average $\widehat{\mu}= n^{-1} \sum_i X_i$:
    \begin{align*}
\mathbb{P}(|\widehat{\mu}-\mu| \geq k \sigma / \sqrt{n}) \leq 2 \exp \left(-k^2 / 2\right).
\end{align*}
\end{cor} 

\begin{remark}
    Actually, we have a tighter Hoeffding bound when the random variables are supported on $[a,b]$ and i.i.d.:
    \[
    \mathbb{P}\left(\left|\hat{\mu}-\mu\right| \geq t/\sqrt{n}\right) \leq 2 \exp \left(-\frac{2 t^2}{(b-a)^2}\right)
    \]
\end{remark}

The following result applies Chernoff bounds to a sequence of sub-Gaussian variables.

\begin{prop}
    [Hoeffding inequality] If $X_i$ are independent $\sigma_i$-subGaussian then for all $t \geq 0$ we have
\begin{align*}
\mathbb{P}\left(\left|\bar{X}_n-\mathbb{E}\left(\bar{X}_n\right)\right| \geq t\right) \leq 2 e^{-t^2 n^2 /\left(2 \sum_i \sigma_i^2\right)}
\end{align*}
\end{prop}


\begin{proof}
    It sufficies to notice that $\bar{X}_n$ is sub-Gaussian with parameter
\begin{align*}
\sigma=\frac{1}{n} \sqrt{\sum_{i=1}^n \sigma_i^2}.
\end{align*}
\end{proof}

\begin{remark}
    If $\sigma_i=\sigma$ for all $i$ we also note that the Hoeffding inequality implies
\begin{align*}
\mathbb{P}\left(\left|\sqrt{n}\left(\bar{X}_n-\mathbb{E}\left(\bar{X}_n\right)\right)\right| \geq t\right) \leq 2 e^{-t^2 /\left(2 \sigma^2\right)}
\end{align*}
\end{remark}

Another general type of bound on the moment generating function is the following.
\begin{de}
    A random variable with mean $\mu=\mathbb{E X}$ is sub-exponential if there are non-negative parameters $v, \alpha$ such that
\begin{align*}
\mathbb{E}\left(e^{\lambda(X-\mu)}\right) \leq e^{v^2 \lambda^2 / 2} \quad \text { for all }|\lambda|<\frac{1}{\alpha} .
\end{align*}
\end{de}

\begin{remark}
    This condition is rather mild, and it is essentially equivalent to the existence of the cumulant generating function in a neighbourhood of zero.
\end{remark}

\begin{eg}
    Let $Z \sim N(0,1)$ and let $X=Z^2$. For $\lambda<1 / 2$, we have
\begin{align*}
\mathbb{E}\left(e^{\lambda(X-1)}\right)=\frac{1}{\sqrt{2 \pi}} \int_{\mathbb{R}} e^{\lambda\left(z^2-1\right)} e^{-z^2 / 2} \mathrm{~d} z= \begin{cases}\frac{e^{-\lambda}}{\sqrt{1-2 \lambda}} & \text { if } \lambda<\frac{1}{2}, \\ +\infty & \text { if } \lambda \geq \frac{1}{2} .\end{cases}
\end{align*}
In particular, since the moment generating function is not defined everywhere, X is not sub-Gaussian. With a bit of calculus, we see that
\begin{align*}
\frac{e^{-\lambda}}{\sqrt{1-2 \lambda}} \leq e^{4 \lambda^2 / 2} \quad \text { for all }|\lambda|<\frac{1}{4}
\end{align*}
and so $X$ is sub-exponential with parameters $(v, \alpha)=(2,4)$.
\end{eg}




With essentially the same proof, the linearity of sub-Gaussian random variables generalizes to subexponential variables.

\begin{lem}
    Suppose that $X=\left(X_1, \ldots, X_n\right)$ is a random vector such that $X_i$ is $\left(v_i, \alpha_i\right)$-sub-exponential for $i=1, \ldots, n$. If $X_i$ are independent and $u \in \mathbb{R}^n$, then the random variable $u^{\top} X$ is sub-exponential with parameters $v=\sqrt{v_1^2 u_1^2+\ldots+v_n^2 u_n^2}$ and $\alpha=\max _i\left|u_i\right| \alpha_i$.
\end{lem}


We obtain simple concentration inequalities for sub-exponential variables.
\begin{prop}
    Suppose that $X$ is sub-exponential with parameters $(v, \alpha)$ and $t \geq 0$. Then
\begin{align*}
\mathbb{P}(X-\mu \geq t) \leq \begin{cases}e^{-\frac{t^2}{2 v^2}} & \text { if } 0 \leq t \leq \frac{v^2}{\alpha} \\ e^{-\frac{t}{2 \alpha}} & \text { if } t>\frac{v^2}{\alpha}\end{cases}
\end{align*}
and
\begin{align*}
\mathbb{P}(|X-\mu| \geq t) \leq \begin{cases}2 e^{-\frac{t^2}{2 v^2}} & \text { if } 0 \leq t \leq \frac{v^2}{\alpha} \\ 2 e^{-\frac{t}{2 \alpha}} & \text { if } t>\frac{v^2}{\alpha}\end{cases}
\end{align*}
\end{prop}

\begin{proof}
    Using the Chernoff bound, we know that
\begin{align*}
\mathbb{P}(X-\mu \geq t) \leq \inf _{\lambda \in\left[0, \frac{1}{\alpha}\right)}\left\{\mathbb{E} e^{\lambda(X-\mathbb{E} X)} e^{-\lambda t}\right\} \leq \inf _{\lambda \in\left[0, \frac{1}{\alpha}\right)}\left\{e^{v^2 \lambda^2 / 2-\lambda t}\right\} .
\end{align*}


The global optimum of $e^{v^2 \lambda^2 / 2-\lambda t}$ is $\lambda^*=t / v^2 \geq 0$. Thus, if $t / v^2< 1 / \alpha$ (equiv. $v^2>t \alpha$ ), we get the sub-Gaussian bound $e^{-t^2 /\left(2 v^2\right)}$.
Otherwise, if $v^2 \leq t \alpha$, the infimum is obtained on the boundary $\lambda^*=1 / \alpha$ and it is equal to $e^{\nu^2 /\left(2 \alpha^2\right)-t / \alpha}$. Note however that the fact that $v^2 \leq t \alpha$, allows to bound this by $e^{-t /(2 \alpha)}$.
\end{proof}

 
\begin{remark}
    These propositions now give a handful of useful results. For example, if $X_i$ are independent sub-exponential with parameters $(v, \alpha)$ then $\bar{X}_n$ is sub-exponential with parameters $\left(\frac{v}{\sqrt{n}}, \frac{\alpha}{n}\right)$ and so
\begin{align*}
\mathbb{P}\left(\left|\bar{X}_n-\mu\right| \geq t\right) \leq \begin{cases}2 e^{-\frac{t^2 n}{2 v^2}} & \text { if } 0 \leq t \leq \frac{v^2}{\alpha} \\ 2 e^{-\frac{t n}{2 \alpha}} & \text { if } t>\frac{v^2}{\alpha}\end{cases}
\end{align*}
\end{remark}



To conclude the introduction to concentration inequalities, we present one of its applications in probabilistic combinatorics.

\begin{de}
    For any integer $n>0$, define the Hamming distance on the hypercube $\{0,1\}^n$
\begin{align*}
d_H(x, y):=\sum_{i=1}^n \mathbf{1}_{x_i \neq y_i}, \quad \text { for any } x, y \in\{0,1\}^n
\end{align*}
\end{de}

\begin{thm} (Gilbert–Varshamov bound for Hamming packing) There exists a universal constant $c>0$, such that for any $n>0$, there exists a subset $A \subseteq\{0,1\}^n$ with $|A| \geq e^{c n}$, satisfying
\begin{align*}
d_H(x, y) \geq \frac{n}{4}, \quad \text { for any pair } x, y \in A .
\end{align*}
\end{thm}

\begin{proof}
    We use a brute-force search. Let $M$ be some constant to be determined, and $X_1,\cdots,X_M$ be i.i.d. samples from the uniform distribution on $\{0,1\}^n$. For any $i\neq j$, since $d_H(X_i,X_j)$ is the sum of $n$ independent Rademacher random variables, we have
    \[
    \pr\paren*{d_H(X_i,X_j)\le n/4}\le 2\exp\paren*{-n/8}.
    \]
    Therefore,
    \[
    \pr\paren*{d_H(x, y) \geq \frac{n}{4},\text { for any pair } x, y \in A }\le M^2  \pr\paren*{d_H(X_i,X_j)\le n/4}\le 2M^2\exp\paren*{-n/8}
    \]
    and it suffices to take $M=\floor{\frac{1}{2}\exp(n/16)}$. Hence, $c=1/16$ suffices.
\end{proof}

\section{Lecture 2, 15: Criteria for Decision Rules}


Now we turn to the framework of decision theory. The goal of classical decision theory is to compare different statistical models and to find or prove the optimality of a specific model.

We start with formulating a statistical model and the method we use to evaluate it.

\begin{de}
A statistical model is a collection of probability distributions $\mathcal{P}=\left\{\mathbb{P}_\theta: \theta \in \Theta\right\}$ of the observed data $X$, parametrized smoothly by the parameter $\Theta$.

Suppose that we observe the data $X \sim \mathbb{P}_{\theta^*}$ for some $\theta^* \in \Theta$, and we have the following attributes of a statistical task:
\begin{enumerate}
    \item A decision rule $\delta$ that maps data $X$ to an action $a \in \mathcal{A}$.
    \item A loss function $L:\Theta\times\mathcal{A}\mapsto \mathbb{R}_+$ that determines the loss of an action under this model.
    \item The risk of a decision rule $\delta$ under the model $\theta$ defined as
    \[
    R(\theta;\delta) =\mathbb{E}_\theta[L(\theta ; \delta(X))] .
    \]
    where $\mathbb{E}_{\theta}$ denotes the expectation under $\pr_{\Theta}$.
\end{enumerate}
\end{de}

We then illustrate some examples of statistical tasks.

\begin{eg}
    (Estimation) The goal is to estimate some quantity of a model, specified as $g(\theta^*)$. To this end, the decision rule $\delta$ is now an estimator $\hat g$ that operates on data, and usually one uses the mean square error to quantify the error of this estimation. Thus, in this task, $L(\theta;\hat g(X))=(g(\theta)-\hat g(X))^2$.
\end{eg}

\begin{eg}
    (Testing) The goal is to decide whether the true parameter $\theta^*$ lies in some subset $\Theta_0$ of probability models. To formulate this, let the action set be $\mathcal{A}=\{0,1\}$, where 0 denotes acceptance and 1 denotes rejection. The loss function is then
    \[
    L(\theta;a)=\begin{cases}
        1,\quad \theta\in\Theta_0, a=1 \text{ or } \theta\notin \Theta_0, a=0,\\
        0,\quad \text{otherwise}.
    \end{cases}
    \]
    As a consequence, $R(\theta;\delta)$ is then the probability of the occurrence of type-I error if $\theta\in\Theta_0$ and the probability of the occurrence of type-II error otherwise.
\end{eg}

In this lecture, we will go over three criteria for comparison of decision rules, including the admissibility of decision rules, Bayes procedures, and minimax procedures, and demonstrate some of their properties.

We start with the definition of admissibility, which is simply the decision rules that are not dominated by any other decision rules under any $\theta\in\Theta$.

\begin{de}
    Suppose $X$ has distribution from a family $\mathcal{P}=\left\{P_\theta: \theta \in \Omega\right\}$ and that $\delta$ is a decision rule. Then $\delta$ is called inadmissible if there exists another decision rule $\tilde{\delta}$ such that
    \[
    R(\tilde\delta;\theta)\le R(\delta;\theta),\quad \forall\theta\in\Theta
    \]
    and strict inequality holds for some $\theta\in\Theta$. If $\delta$ is not inadmissible, then $\delta$ is admissible.
\end{de}

The problem with this straightforward definition is that, sometimes, naive decision rules may be admissible and classical decision rules may be inadmissible. See the following examples.

\begin{eg}
    (Naive decision rules may be admissible) Suppose we have i.i.d. random variables $X_1, X_2, \ldots, X_n \sim \mathcal{N}\left(\theta, \sigma^2 \right)$ where $\theta\in \mathbb{R}$ is unknown and $\sigma^2$ is known. Our goal is to estimate $\theta$ with an estimator $\hat{\theta}$ under the mean-square loss. 

    We now show that the naive estimator $\hat{\delta}\equiv0$ is admissible. Indeed, if some $\tilde\delta$ strictly dominates $\hat{\delta}$, we know that $\tilde\delta$ will achieve zero loss in the case that $\theta=0$. This indicates that $\tilde\delta(X)=0$ for all possible $X=(X_1, X_2, \ldots, X_n)$. Since the Gaussian distribution has full support over the real line, we know that $\tilde\delta$ is identically zero and hence $\hat{\delta}$ is admissible.
\end{eg}


Based on this drawback, it is reasonable to find some features that capture the risk $R(\delta;\theta)$ for all $\theta$ instead of considering them simultaneously. There are two main approaches: the first one is to consider a weighted average of the loss function among $\theta$, which leads to a Bayes procedure; the other one is to consider the worst case among all $\theta$, which leads to a minimax procedure.

\begin{de}
     Suppose $X$ has distribution from a family $\mathcal{P}=\left\{P_\theta: \theta \in \Omega\right\}$ and that we impose a prior $\pi$ on $\Theta$. The Bayes risk is then defined as
     \[
     r_\pi(\delta)=\int_\Theta \E_\theta\brackets*{L(\theta;\delta(X))} \pi(\d \theta),
     \]
     and the Bayes decision rule is then
     \[
     \delta_{\operatorname{Bayes},\pi}=\operatornamewithlimits{argmin}_{\delta}  r_\pi(\theta)
     \]
\end{de}

\begin{prop}
    Suppose that the distribution of $X\sim P_\theta$ has density $p_\theta$ with respect to some base measure $\lambda$. Then the Bayes decision rule is given by
    $$\delta_{\text{Bayes},\pi}(\mathbf{x}) = \underset{a \in A}{\operatorname{argmin}} \E_{\pi(\cdot|\mathbf{x})}[L(\theta,a)],$$
where $\pi(\theta|\mathbf{x}) = \frac{p_{\theta}(\mathbf{x})}{\int_{\Theta} p_{\theta'}(\mathbf{x}) \pi(d\theta')}\pi(d\theta) $ is the posterior distribution of $\theta$.
\end{prop}

\begin{proof}
    Since the density of $X$ is given by $p_\theta$, we have
    $$ 
    r_{\pi}(\delta) = \int_{\Theta} \int_{X} L(\theta; \delta(\mathbf{x})) p_{\theta}(x) \lambda(d\mathbf{x})\pi(d\theta) = \int_{\mathbf{X}} \left[ \int_{\Theta} L(\theta; \delta(\mathbf{x})) p_{\theta}(\mathbf{x}) \pi(d\theta) \right] \lambda(d\mathbf{x}),
    $$
and the result can be seen from that for each fixed $\mathbf{x}$, $\pi(\theta|\mathbf{x})\propto p_\theta(\mathbf{x})\pi(d\theta)$.
\end{proof}

\begin{de}
     Suppose $X$ has distribution from a family $\mathcal{P}=\left\{P_\theta: \theta \in \Omega\right\}$. The minimax decision rule is defined as
     \[
     \delta_{\operatorname{minimax}}=\operatornamewithlimits{argmin}_{\delta}  \sup_{\theta} R(\theta;\delta)
     \]
\end{de}

\begin{remark}
    From a game-theoretic perspective, $\delta_{\operatorname{minimax}}$ can be viewed as the optimal strategy of the zero-sum game between the statistician and the environment with the reward $R(\theta;\delta)$. As a result, we have the following propositions from game theory.
\end{remark}

\begin{prop}
    (Weak duality) Suppose $X$ has distribution from a family $\mathcal{P}=\left\{P_\theta: \theta \in \Omega\right\}$. Then it holds that
    $$
    \sup_\theta \inf_\delta R(\theta;\delta) \leq \inf_\delta \sup_\theta  R(\theta;\delta)
    $$
\end{prop}

\begin{prop}
    (Strong duality) Suppose in addition that the spaces $\Theta$, $\mathcal{A}$, and $X$ are finite, then strong duality holds:
    $$
    \sup_\theta \inf_\delta R(\theta;\delta) = \inf_\delta \sup_\theta  R(\theta;\delta)
    $$
\end{prop}

Now we turn to relations and properties between these procedures.

\begin{prop}
     Bayes rules with constant risks are minimax. 
\end{prop}

\begin{proof}
Let $\theta$ be a Bayes rule under prior $\pi$ with constant risk. Consider another decision rule $\delta^{\prime}$, then
\begin{align*}
\begin{aligned}
\sup _{\theta \in \Theta} R\left(\theta; \delta^{\prime}\right) & \geq \int_{\Theta} R\left(\theta, \delta^{\prime}\right) \pi(d \theta)\geq \int_{\Theta} R\left(\theta, \delta\right) \pi(d \theta)=\sup _{\theta \in \Theta} R(\theta, \delta)
\end{aligned}
\end{align*}
\end{proof}

The next result gives a connection between minimax and Bayesian rules as a generalization of the previous results.
\begin{thm}
    Let $\delta^*$ be a Bayes rule for some prior $\pi$. Suppose that
\begin{align*}
R\left(\theta, \delta^*\right) \leq r\left(\pi, \delta^*\right) \quad \text { for all } \quad \theta \in \Theta .
\end{align*}
Then $\delta^*$ is minimax.
\end{thm}

\begin{proof}
    We prove the contrapositive statement. If $\delta^*$ is not minimax then there exists $\delta_0$ such that $\bar{R}\left(\delta_0\right)<\bar{R}\left(\delta^*\right)$. However, the fact that $r\left(\pi, \delta_0\right) \leq \sup _\theta R\left(\theta, \delta_0\right)$ implies that
\begin{align*}
r\left(\pi, \delta_0\right) \leq \bar{R}\left(\delta_0\right)<\bar{R}\left(\delta^*\right) \leq r\left(\pi, \delta^*\right) .
\end{align*}
But this contradicts the fact that $\delta^*$ is a Bayesian rule.
\end{proof}

\begin{prop}
    Suppose that $\exists\left(\pi_j, r_j\right)_{j=1}^{\infty}$ such that $r_j$ is the Bayes risk under $\pi_j$ using the Bayes decision rule. Suppose that there exists another decision rule $\delta$ such that $\liminf _{j\rightarrow+\infty} r_j \geq \sup _\theta R(\theta, \delta)$, then $\delta$ is minimax.
\end{prop}

\begin{proof}
    For another decision rule $\delta^{\prime}$, we have
\begin{align*}
\begin{aligned}
\sup _\theta R\left(\theta, \delta^{\prime}\right) & \geq\liminf _{j \rightarrow+\infty} \int_{\Theta} R\left(\theta, \delta^{\prime}\right) \pi_j(d \theta) \\
& \geq \liminf _{j \rightarrow+\infty} \int_{\Theta} R\left(\theta, \delta_{\text {Bayes }}\right) \pi_j(d \theta) \\
& =\liminf _{j \rightarrow+\infty} r_j \geq \sup_\theta R(\theta, \delta) .
\end{aligned}
\end{align*}
\end{proof}

\begin{prop}
Unique Bayes rules $\delta$ are admissible.
\end{prop}

\begin{proof}
    Suppose $\delta^{\prime}$ dominates $\delta$, then we have
\begin{align*}
r_\pi\left(\delta^{\prime}\right)=\int R\left(\theta; \delta^{\prime}\right) \pi(d \theta) \leq \int R(\theta; \delta) \pi(d \theta)=r_\pi(\delta)
\end{align*}

So $\delta^{\prime}=\delta$ by uniqueness.
\end{proof}

\begin{prop}
    When $|\Theta|<+\infty$, admissible rules are Bayes (under some prior $\pi$).
\end{prop}

\begin{proof}
Let $\delta_0$ be admissible and $|\Theta|=k$. Consider the convex set
\[
T := \{\,\left(R(\theta_i;\delta)-R(\theta_i;\delta_0)\right)_{i=1,\cdots,k}:\ \delta\ \text{a rule}\,\}\subset\mathbb{R}^k.
\]
Because $R(\theta;\cdot)$ is convex (mixtures of randomized rules produce convex combinations of risk vectors), $T$ is convex and contains the origin $0$ (corresponding to $\delta=\delta_0$). Admissibility of $\delta_0$ means precisely that there is no $t\in T$ with $t\le 0$ (component-wise) and $t\neq 0$. Equivalently,
\[
T\cap\big(-\mathbb{R}^k_{++}\big)=\varnothing,
\]
where $\mathbb{R}^k_{++}=\{x\in\mathbb{R}^k:\ x_i>0,\ \forall i\}$.

The sets $T$ and $-\,\mathbb{R}^k_{++}$ are convex and disjoint, and $-\,\mathbb{R}^k_{++}$ is an open convex cone. By the separating hyperplane theorem, there exists a nonzero vector $\pi\in\mathbb{R}^k$ and a scalar $c$ such that
\[
\pi\cdot t \ge c \quad\text{for all } t\in T,
\qquad\text{and}\qquad
\pi\cdot y \le c \quad\text{for all } y\in -\mathbb{R}^k_{++}.
\]
Because $0\in T$, plugging $t=0$ gives $0=\pi\cdot 0 \ge c$, so $c\le0$. Replace $\pi$ by $-\pi$ if necessary so that $c\le 0$ still holds; we may then take $c=0$ by scaling the separating hyperplane (the sign and scale of $\pi$ are not important for separation). Thus we obtain a nonzero $\pi$ with
\[
\pi\cdot t \ge 0 \quad\text{for all } t\in T,
\]
and
\[
\pi\cdot y \le 0 \quad\text{for all } y\in -\mathbb{R}^k_{++}.
\]

The second condition implies that $\pi_i\ge 0$ for every $i$: indeed, taking $y=-e_i\in -\mathbb{R}^k_{++}$ (where $e_i$ is the $i$th standard basis vector) yields $-\pi_i\le 0$, hence $\pi_i\ge0$. Not all $\pi_i$ are zero because $\pi\neq 0$. Normalize $\pi$ by dividing by $\sum_i\pi_i$ to obtain a probability vector (prior) which we again denote by $\pi=(\pi_1,\dots,\pi_k)$ with $\pi_i\ge0$, $\sum_i\pi_i=1$.

Now for any decision rule $\delta$ let $t=R(\delta)-R(\delta_0)\in T$. The inequality $\pi\cdot t\ge0$ becomes
\[
\sum_{i=1}^k \pi_i \big(R(\theta_i,\delta)-R(\theta_i,\delta_0)\big)\ge 0,
\]
or equivalently
\[
r_\pi(\delta)\ge r_\pi(\delta_0).
\]
Therefore $\delta_0$ attains the minimal Bayes risk with respect to the prior $\pi$, i.e. $\delta_0$ is a Bayes rule for $\pi$. This completes the proof.
\end{proof}




\section{Lecture 3, 16: (Unbiased) Estimation}
\subsection{Construction of UMVUE}

An estimator is called unbiased if $\mathbb{E}_\theta \delta(X)=\theta$ for all $\theta$. Sometimes the following definition is used instead.
\begin{de}
    For a loss $L(\theta, a)$ a decision rule $\delta$ is unbiased with respect to $L$ if
\begin{align*}
\mathbb{E}_\theta\left(L\left(\theta^{\prime}, \delta(X)\right)\right) \geq \mathbb{E}_\theta(L(\theta, \delta(X)))=R(\theta, \delta) \quad \text { for all } \theta, \theta^{\prime} \in \Theta .
\end{align*}
\end{de}



We start with the introduction of sufficient statistics.
\begin{eg}
    Suppose that we have observations $X_1, X_2 \ldots X_n \stackrel{\text { iid }}{\sim} \operatorname{Ber}(p)$ from a Bernoulli distribution with an unknown parameter $p$. Then we know that the statistics $T\left(X_1,\cdots,X_n\right)=X_1+\cdots+X_n \sim \operatorname{Binom}(n, p)$ follows a Binomial distribution. 
    
    Moreover, based on this statistics, we can obtain the distribution of the data without knowing the exact value of $p$: the conditional distribution of $X=(X_1,\cdots,X_n)$ conditioned on $\{T=t\}$ is the uniform distribution among all subsets with $t$ ones in $\{0,1\}^n$.
\end{eg}
This idea can be further generalized to other cases.

 \begin{de}
     A statistics $T$ is called sufficient if the distribution of the data $X$ is independent of $\theta$ given $T(X)$. In other words, there exists a distribution $\mu_X$ such that
     \[
     \mu_X(B)=\pr_\theta\left(X\in B|T(X)=t\right),\quad \forall\theta\in\Theta, B \in \mathcal{B}, t\in \mathbb{R}.
     \]
 \end{de}

We have the following nice characterization of sufficient statistics.
\begin{de}
    A family of distributions $\mathcal{P}=\left\{P_\theta: \theta \in \Omega\right\}$ is dominated if there exists a measure $\mu$ with $P_\theta$ absolutely continuous with respect to $\mu$, for all $\theta \in \Omega$.
\end{de}

\begin{thm}
    (Factorization Theorem). Let $\mathcal{P}=\left\{P_\theta: \theta \in \Omega\right\}$ be a family of distributions dominated by $\mu$. A necessary and sufficient condition for a statistic $T$ to be sufficient is that there exist functions $g_\theta \geq 0$ and $h \geq 0$ such that the densities $p_\theta$ for the family satisfy
\begin{align*}
p_\theta(x)=g_\theta(T(x)) h(x), \text { for a.e. } x \text { under } \mu .
\end{align*}
\end{thm}

\begin{proof}
    Necessity: We construct $g_\theta$ and $h$ as
\begin{align*}
g_\theta(t)=P_\theta(T(x)=t), \quad h(x)=P_{\theta_0}(X=x \mid T(x)=t)
\end{align*}
Then by definition of sufficiency, $h(x)$ is independent of $t$. Therefore,
\begin{align*}
P_\theta(x)=P_\theta(X=x)=P_\theta(T(x)=t) P_\theta\left(X=x(T(x)=t)=g_\theta(T(x)) h(x) .\right.
\end{align*}


Sufficiency: Assume $P_\theta(x)=\rho_\theta(T(x)) h(x)$. Then $\forall \theta_0 \in \mathbb{H}, t=T(x)$
\begin{align*}
\begin{aligned}
P_{\theta_0}(X=x \mid T(X)=t) & =\frac{P_{\theta_0}(X=x, T(X)=t)}{P_{\theta_0}(T(X)=t)}=\frac{P_{\theta_0}(X=x)}{\int_{T(x)=t} P_{\theta_0}(x) d x} \\
& =\frac{P_{\theta_0}(x)}{\int_{T(x)=t} P_{\theta_0}(x) d x}=\frac{g_{\theta_0}(T(x)) h(x)}{g_{\theta_0}(t) \int_{T(x)=t} h(x) d x} \\
& =\frac{h(x)}{\int_{T(x)=t} h(x) d x}
\end{aligned}
\end{align*}
which is independent of $\theta_0 \in \Theta$.
\end{proof}

\begin{eg}
    Let $P_\theta(x)=h(x) \exp \left(\eta(\theta)^{\top} T(x)-B(\theta)\right)$ be an exponential family. Then $T(X)$ is sufficient. Moreover, if $X_1, \cdots, X_n \sim P_\theta(x)$. then $\sum_{i=1}^n T\left(X_i\right)$ is sufficient.
\end{eg}

In essence, the complete statistics capture all the information available in the distribution family $\mathcal{P}$. To gather this information with minimal effort, we need the notion of minimal statistics.
\begin{de}
$T(X)$ is a minimal sufficient statistic if
\begin{itemize}
    \item It is a sufficient statistic for $\theta$.
    \item For any other sufficient statistic $S(X)$, there exists a function $g$ such that $T(X)= g(S(X))$.
\end{itemize}
\end{de}

This gives rise to the so-called complete statistics.

\begin{de}
    A statistic $T$ is complete for a family $\mathcal{P}=\left\{P_\theta: \theta \in \Omega\right\}$ if
\begin{align*}
E_\theta f(T)=0, \text { for all } \theta,
\end{align*}
imply $f(T)=0$ (a.e. $\mathcal{P}$ ).
\end{de}


\begin{eg}
    Suppose $X_1, \ldots, X_n$ are i.i.d. from a uniform distribution on $(0, \theta)$. Using indicator functions, the joint density is $I\left\{\min x_i>0\right\} I\left\{\max x_i<\right. \theta\} / \theta^n$, and so $T=\max \left\{X_1, \ldots, X_n\right\}$ is sufficient by the factorization theorem. Note that $T$ has density $n t^{n-1} / \theta^n, t \in(0, \theta)$. Suppose $E_\theta f(T)=0$ for all $\theta>0$; then
\begin{align*}
E_\theta[f(T)]=\frac{n}{\theta^n} \int_0^\theta[f(t)] t^{n-1} d t=0
\end{align*}
From this, $[f(t)-c] t^{n-1}=0$ for a.e. $t>0$. So $f(T)=c$ (a.e. $\mathcal{P}$ ), and $T$ is complete.
\end{eg}


\begin{thm}
    Suppose $\mathcal{P}=\left\{P_\theta: \theta \in \Omega\right\}$ is a dominated family with densities $p_\theta(x)=g_\theta(T(x)) h(x)$. If $p_\theta(x) \propto{ }_\theta p_\theta(y)$ implies $T(x)=T(y)$, then $T$ is minimal sufficient.
\end{thm}

\begin{proof}
    A proper proof of this result unfortunately involves measure-theoretic niceties, but here is the basic idea. Suppose $\tilde{T}$ is sufficient. Then $p_\theta(x)= \tilde{g}_\theta(\tilde{T}(x)) \tilde{h}(x)$ (a.e. $\mu$ ). Assume this equation holds for all $x$. If $T$ is not a function of $\tilde{T}$, then there must be two data sets $x$ and $y$ that give the same value for $\tilde{T}, \tilde{T}(x)=\tilde{T}(y)$, but different values for $T, T(x) \neq T(y)$. But then
\begin{align*}
p_\theta(x)=\tilde{g}_\theta(\tilde{T}(x)) \tilde{h}(x) \propto_\theta \tilde{g}_\theta(\tilde{T}(y)) \tilde{h}(y)=p_\theta(y),
\end{align*}
and from the condition on $T$ in the theorem, $T(x)$ must equal $T(y)$. Thus $T$ is a function of $\tilde{T}$. Because $\tilde{T}$ was an arbitrary sufficient statistic, $T$ is minimal.
\end{proof}

\begin{thm}
    If $T$ is complete and sufficient, then $T$ is minimal sufficient.
\end{thm}
\begin{proof}
    Let $\tilde{T}$ be a minimal sufficient statistic, and assume $T$ and $\tilde{T}$ are both bounded random variables. Then $\tilde{T}=f(T)$. Define $g(\tilde{T})=E_\theta[T \mid \tilde{T}]$, noting that this function is independent of $\theta$ because $\tilde{T}$ is sufficient. By smoothing, $E_\theta g(\tilde{T})=E_\theta T$, and so $E_\theta[T-g(\tilde{T})]=0$, for all $\theta$. But $T-g(\tilde{T})=T-g(f(T))$, a function of $T$, and so by completeness, $T=g(\tilde{T})$ (a.e. $\mathcal{P}$ ). This establishes a one-to-one relationship between $T$ and $\tilde{T}$. From the definition of minimal sufficiency, $T$ must also be minimal sufficient.

    For the general case, first note that sufficiency and completeness are both preserved by one-to-one transformations, so two statistics can be considered equivalent if they are related by a one-to-one (bimeasurable) function. But there are one-to-one bi-measurable functions from $\mathbb{R}^n$ to $\mathbb{R}$, and so any random vector is equivalent to a single random variable. Using this, if $T$ and $\tilde{T}$ are random vectors, the result follows easily from the one-dimensional case, transforming both of them to equivalent random variables.
\end{proof}

We now extend the previous example and show that the statistics $T$ of a (sufficiently regular) exponential family are minimal.

\begin{de}
    An exponential family with densities $p_\theta(x)=\exp \{\eta(\theta) \cdot T(x)-B(\theta)\} h(x), \theta \in \Omega$, is said to be of full rank if the interior of $\eta(\Omega)$ is not empty and if $T_1, \ldots, T_s$ do not satisfy a linear constraint of the form $v \cdot T=c$ (a.e. $\mu$).
\end{de}

\begin{remark}
    If $\Omega \subset \mathbb{R}^s$ and $\eta$ is continuous and one-to-one (injective), and the interior of $\Omega$ is nonempty, then the interior of $\eta(\Omega)$ cannot be empty. This follows from the "invariance of domain" theorem of Brouwer.
\end{remark}


If the interior of $\eta(\Omega)$ is not empty, then the linear span of $\eta(\Omega) \ominus \eta(\Omega)$ will be all of $\mathbb{R}^s$, and, $T$ will be minimal sufficient. The following result shows that in this case $T$ is also complete.

\begin{thm}
    In an exponential family of full rank, $T$ is complete.
\end{thm}

\begin{proof}
    Let $\eta\left(\theta_0\right)$ be an interior point of $\eta(\Theta)$. Suppose $\mathbb{E}_\theta[f(T(x))]=0$ for some $f$, then $\forall \theta \in \Theta$,
\begin{align*}
\begin{aligned}
0 & =\int_{\mathbb{R}^d} f(t) h(t) \exp \left(\eta^{\top}(\theta) t-A(\theta)\right) d t \\
\Rightarrow 0 & =\int_{\mathbb{R}^d} f(t) h(t) \exp \left(\eta^{\top}(\theta) t\right) d t
\end{aligned}
\end{align*}
So the Laplace transform of $f \cdot v$ is 0, by the uniqueness of the Laplace transform $f(t) v(t)\equiv0$ and hence $f \equiv 0$.
\end{proof}

\begin{de}
    A statistic $V$ is called ancillary if its distribution does not depend on $\theta$. 
\end{de}

An ancillary statistic $V$, by itself, provides no useful information about $\theta$. But in some situations $V$ can be a function of a minimal sufficient statistic $T$. The following result of Basu shows that when $T$ is complete it will contain no ancillary information. 
\begin{thm}[Basu]
 If $T$ is complete and sufficient for $\mathcal{P}=\left\{P_\theta: \theta \in\right. \Omega\}$, and if $V$ is ancillary, then $T$ and $V$ are independent under $P_\theta$ for any $\theta \in \Omega$.
\end{thm}

\begin{proof}
    Define $q_A(t)=P_\theta(V \in A \mid T=t)$, so that $q_A(T)=P_\theta(V \in A \mid T)$, and define $p_A=P_\theta(V \in A)$. By sufficiency and ancillarity, neither $p_A$ nor $q_A(t)$ depend on $\theta$. Also, by smoothing,
\begin{align*}
p_A=P_\theta(V \in A)=E_\theta P_\theta(V \in A \mid T)=E_\theta q_A(T),
\end{align*}
and so, by completeness, $q_A(T)=p_A$ (a.e. $\mathcal{P}$ ). By smoothing,
\begin{align*}
\begin{aligned}
P_\theta(T \in B, V \in A) & =E_\theta 1_B(T) 1_A(V) \\
& =E_\theta E_\theta\left(1_B(T) 1_A(V) \mid T\right) \\
& =E_\theta 1_B(T) E_\theta\left(1_A(V) \mid T\right) \\
& =E_\theta 1_B(T) q_A(T) \\
& =E_\theta 1_B(T) p_A \\
& =P_\theta(T \in B) P_\theta(V \in A) .
\end{aligned}
\end{align*}
Here $A$ and $B$ are arbitrary Borel sets, and so $T$ and $V$ are independent.
\end{proof}


\begin{eg}
    Suppose $X_1, \ldots, X_n$ are i.i.d. from $N\left(\mu, \sigma^2\right)$, and take $\mathcal{P}= \mathcal{P}_\sigma=\left\{N\left(\mu, \sigma^2\right)^n: \mu \in \mathbb{R}\right\}$. (Thus $\mathcal{P}_\sigma$ is the family of all normal distributions with standard deviation the fixed value $\sigma$.) With $\bar{x}=\left(x_1+\cdots+x_n\right) / n$, the joint density can be written as
\begin{align*}
\frac{1}{\left(2 \pi \sigma^2\right)^{n / 2}} \exp \left[\frac{n \mu}{\sigma^2} \bar{x}-\frac{n \mu^2}{2 \sigma^2}-\frac{1}{2 \sigma^2} \sum_{i=1}^n x_i^2\right] .
\end{align*}


These densities for $\mathcal{P}_\sigma$ form a full rank exponential family, and so the average $\bar{X}=\left(X_1+\cdots+X_n\right) / n$ is a complete sufficient statistic for $\mathcal{P}_\sigma$. Define
\begin{align*}
S^2=\frac{1}{n-1} \sum_{i=1}^n\left(X_i-\bar{X}\right)^2
\end{align*}
called the sample variance. For the family $\mathcal{P}_\sigma, S^2$ is ancillary. To see this, let $Y_i=X_i-\mu, i=1, \ldots, n$. Because
\begin{align*}
\begin{aligned}
P_\mu\left(Y_i \leq y\right)=P_\mu\left(X_i \leq y+\mu\right) & =\int_{-\infty}^{y+\mu} \exp \left[-\frac{1}{2 \sigma^2}(x-\mu)^2\right] \frac{d x}{\sqrt{2 \pi \sigma^2}} \\
& =\int_{-\infty}^y \exp \left[-\frac{1}{2 \sigma^2} u^2\right] \frac{d u}{\sqrt{2 \pi \sigma^2}}
\end{aligned}
\end{align*}
and the integrand is the density for $N\left(0, \sigma^2\right), Y_i \sim N\left(0, \sigma^2\right)$. Then $Y_1, \ldots, Y_n$ are i.i.d. from $N\left(0, \sigma^2\right)$. Because $\bar{Y}=\left(Y_1+\cdots+Y_n\right) / n=\bar{X}-\mu, X_i-\bar{X}= Y_i-\bar{Y}, i=1, \ldots, n$, and
\begin{align*}
S^2=\frac{1}{n-1} \sum_{i=1}^n\left(Y_i-\bar{Y}\right)^2
\end{align*}
Because the joint distribution of $Y_1, \ldots, Y_n$ depends on $\sigma$ but not $\mu, S^2$ is ancillary for $\mathcal{P}_\sigma$. Hence, by Basu's theorem, $\bar{X}$ and $S^2$ are independent.
\end{eg}




 The existence of sufficient statistics allows us to operate on the statistics, which have a fixed distribution based on $T$, instead of the data itself, which may have different distributions as $\theta$ varies.

 \begin{thm}
     Suppose $X$ has distribution from a family $\mathcal{P}=\left\{P_\theta: \theta \in \Omega\right\}$ and that $T=T(X)$ is sufficient. Then, for any decision rule $\delta(X)$ of $g(\theta)$, there exists a randomized decision rule $\tilde{\delta}$ based on $T$ that has the same distribution as $T(X)$. Consequently, it has the same risk function as $\delta(X)$.
 \end{thm}

\begin{proof}
    It suffices to take $\tilde\delta(T(X))=\delta(Y)$, where $Y$ is distributed as $X|T(X)$.
\end{proof}

 \begin{thm}
     (Rao-Blackwellization) Suppose $X$ has distribution from a family $\mathcal{P}=\left\{P_\theta: \theta \in \Omega\right\}$ and that $T=T(X)$ is sufficient. If $L(\theta, \cdot)$ is convex in the action space $\mathcal{A}$, then for any decision rule $\delta$, the new decision rule $\eta(T(X)):=\mathbb{E}[\delta(X) \mid T(x)]$ that operates on $T(X)$ satisfies $R(\theta ; \eta) \leq R(\theta ; \delta)$. 
 \end{thm}

\begin{proof}
    Follows from an application of Jensen's inequality.
\end{proof}

\begin{remark}
    The estimated risk $R(\theta ; \eta)$ is computable since it does not rely on the model parameter $\theta$.
\end{remark}

\begin{de}
    An estimator $\delta$ is called unbiased for $g(\theta)$ if
\begin{align*}
E_\theta \delta(X)=g(\theta), \quad \forall \theta \in \Omega .
\end{align*}
If an unbiased estimator exists, $g$ is called U-estimable.
\end{de}



\begin{de}
    An unbiased estimator $\delta$ is uniformly minimum variance unbiased (UMVU) if
\begin{align*}
\operatorname{Var}_\theta(\delta) \leq \operatorname{Var}_\theta\left(\delta^*\right), \quad \forall \theta \in \Omega,
\end{align*}
for any competing unbiased estimator $\delta^*$.
\end{de}

Since the space of unbiased estimators is a linear space and UMVU acts as a projection under the $L^2$ norm, we have the following theorem:
\begin{thm}
An unbiased estimator $T$ is the UMVUE for its expected value if and only if its covariance with any unbiased estimator of zero is exactly zero. That is, $\operatorname{Cov}(T, U)=0$ for all statistics $U$ where $E[U]=0$.
\end{thm}

In a general setting, there is no reason to suspect that there will be a UMVU estimator. However, if the family has a complete sufficient statistic, a UMVU will exist, at least when $g$ is U-estimable.

\begin{thm}
    Suppose $g$ is $U$-estimable and $T$ is complete sufficient. Then there is an essentially unique unbiased estimator based on $T$ that is UMVU.
\end{thm}
\begin{proof}
    Let $\delta=\delta(X)$ be any unbiased estimator and define
\begin{align*}
\eta(T)=E[\delta \mid T],
\end{align*}
as in the Rao-Blackwell theorem. By smoothing,
\begin{align*}
g(\theta)=E_\theta \delta=E_\theta E_\theta[\delta \mid T]=E_\theta \eta(T),
\end{align*}
and thus $\eta(T)$ is unbiased. Suppose $\eta^*(T)$ is also unbiased. Then
\begin{align*}
E_\theta\left[\eta(T)-\eta^*(T)\right]=0, \quad \forall \theta \in \Omega,
\end{align*}
and by completeness, $\eta(T)-\eta^*(T)=0$ (a.e. $\mathcal{P}$ ). This shows that the estimator $\eta(T)$ is essentially unique; any other unbiased estimator based on $T$ will equal $\eta(T)$ except on a $\mathcal{P}$-null set. The estimator $\eta(T)$ has minimum variance by the Rao-Blackwell theorem with squared error loss. Specifically, if $\delta^*$ is any unbiased estimator, then $\eta^*(T)=E_\theta\left(\delta^* \mid T\right)$ is unbiased by the calculation above. With squared error loss, the risk of $\delta^*$ or $\eta^*(T)$ is the variance, and so
\begin{align*}
\operatorname{Var}_\theta\left(\delta^*\right) \geq \operatorname{Var}_\theta\left(\eta^*(T)\right)=\operatorname{Var}_\theta(\eta(T)), \quad \forall \theta \in \Omega
\end{align*}
Thus, $\eta(T)$ is UMVU.
\end{proof}

\begin{remark}
    This also gives a way to compute the UMVU.
\end{remark}

\subsection{Lower Bounds of Risks of Unbiased Estimators}

We now introduce the famous Cramer-Rao lower bound of unbiased estimators. We start with the definition of properties of score functions.

\begin{lem}[Fisher's identity]
    Let $l(\theta;X) =  \log p_\theta(X)$. Then
\begin{align*}
\begin{aligned}
\mathbb{E}_\theta\left[\nabla_\theta l(\theta ; X)\right] & =\int\nabla_\theta l(\theta ; x)p_\theta(x) d \mu(x)=\nabla_\theta\left(\int p_\theta(x) d \mu(x)\right)=0 .
\end{aligned}
\end{align*}
\end{lem}

\begin{lem}
    From Fisher's identity, we have
\begin{align*}
\mathbb{E}_\theta\left[\frac{\partial}{\partial \theta_j} l(\theta ; X)\right]=0 .
\end{align*}
Therefore,
\begin{align*}
\begin{aligned}
0 & =\frac{\partial}{\partial \theta_i} \mathbb{E}_\theta\left[\frac{\partial}{\partial \theta_j} l(\theta ; x)\right] =\frac{\partial}{\partial \theta_i} \int \frac{\partial}{\partial \theta_j} l(\theta ; x) \cdot e^{l(\theta ; x)} d \mu(x)  \\
& =\int\left(\frac{\partial^2 l(\theta ; x)}{\partial \theta_j \partial \theta_i}+\frac{\partial l(\theta ; x)}{\partial \theta_i} \cdot \frac{\partial l(\theta, x)}{\partial \theta_j}\right) e^{l(\theta ; x)} d \mu(x).
\end{aligned}
\end{align*}

and we have
\begin{align*}
\mathbb{E}_\theta\left[\nabla \ell(\theta ; X) \nabla \ell(\theta ; X)^{\top}\right]=-\mathbb{E}_\theta\left[\nabla^2 \ell(\theta ; X)\right] .
\end{align*}

\end{lem}

\begin{thm}[Cramer-Rao lower bound]    
Let $\delta$ be an unbiased estimator for $g(\theta) \in \mathbb{R}$. Assume that $(p_\theta: \theta \in \Theta)$ has a shared support for all $\theta$, and 
$$
\log p_\theta(x) \in C^2, \quad \mathbb{E}\left|\nabla \log p_\theta(x)\right|^2<+\infty.
$$
Then we have the Cramer-Rao lower bound
\begin{align*}
\var_\theta(\delta(X)) \geq \nabla g(\theta)^{\top} I(\theta)^{-1} \nabla g(\theta)
\end{align*}
where the Fisher information matrix is given by
$$I(\theta):=\mathbb{E}_\theta\left[\nabla \log p_\theta(x) \nabla \log p_\theta(x)^{\top}\right].$$
\end{thm}

\begin{proof}
From Fisher's identity, we have
\begin{align*}
\int \nabla_\theta l(\theta ; x) e^{l (\theta ; x)} d \mu(x)=0 .
\end{align*}
and therefore,
\begin{align*}
\begin{aligned}
\nabla g(\theta) & =\int(\delta(x)-g(\theta)) \cdot \nabla_\theta l(\theta ; x) \cdot e^{l(\theta ; x)} d \mu(x) \\
& =\mathbb{E}_\theta\left[(\delta(x)-g(\theta)) \cdot \nabla_\theta l(\theta ; x)\right] .
\end{aligned}
\end{align*}
Consider the quadratic form
\begin{align*}
u \in \mathbb{R}^d \mapsto \mathbb{E}_0\left[\left(u^{\top} \nabla_\theta l(\theta ; x)-(\delta(x)-g(\theta))\right)^2\right] 
\end{align*}
which is nonnegative, a direct expansion with respect to $u$ suggests that
\begin{align*}
0 \leq u^{\top} I(\theta) u-2 u^{\top} \mathbb{E}_\theta\left[(\delta(x)-g(\theta)) \cdot \nabla_\theta l(\theta ; x)\right]+\var_\theta[\delta(X)]
\end{align*}
The conclusion then follows from noting that $\mathbb{E}_\theta\left[(\delta(x)-g(\theta)) \cdot \nabla_\theta l(\theta ; x)\right]=\nabla g(\theta)$.
\end{proof}

\begin{cor}
Let $g: \mathbb{R}^d \rightarrow \mathbb{R}^k$ and $\delta$ be a multi-dimensional unbiased estimator of $g(\theta)$. Then
\begin{align*}
\mathbb{E}\left[(\delta(X)-g(\theta))^{\otimes 2}\right]  \succeq\nabla g(\theta)^{\top} I(\theta)^{-1} \nabla g(\theta).
\end{align*}
In particular, when $g(\theta)=\theta$,
\begin{align*}
\mathbb{E}\left[\|\delta(X)-\theta \|_2^2\right] \geq \operatorname{Tr}\left(I(\theta)^{-1}\right) .
\end{align*}
\end{cor}

\begin{remark}
    The CR lower bound is usually not achieved by unbiased estimation. However, it is achieved asymptotically by the maximum likelihood estimator.
\end{remark}

\begin{remark}
    To attain the CRLB, an estimator must satisfy the condition for equality in the Cauchy-Schwarz inequality used to derive the bound. This occurs if and only if the estimator is a linear function of the score function.

For regular exponential families, the score function is linear in the sufficient statistic $S$. Therefore, an unbiased estimator can only achieve the CRLB if it is a linear function of the sufficient statistic $S$.
\end{remark}

\begin{remark}
    When $p_\theta(x)$ is singular or has discontinuities, the CR lower bound may not hold. For example, consider $ X_1,X_2,\cdots,X_n \sim \operatorname{Unif}([0,\theta])$. Then the UPVU is given by $f(X)=\frac{n+1}{n}\max _{i} X_i$, and $\operatorname{Var}_\theta[f(X)]=\frac{\theta^2}{n(n+2)}$, which is better than the CR lower bound.
\end{remark}


We now introduce the Van-Trees inequality, also known as the Bayesian Cramér-Rao bound, which extends the classical Cramér--Rao inequality to the Bayesian setting. Consider the Bayes risk
\begin{align*}
r_\pi(\delta):=\int_{\Theta} \mathbb{E}_\theta\left[(\delta(x)-g(\theta))^2\right] \pi(d \theta) .
\end{align*}
\begin{thm}[van-Trees inequality]
    For any estimator $\delta$ and prior $\pi$ compactly supported on $\Theta$, we have
\begin{align*}
r_\pi(\delta) \geq\left(\int_{\Theta} \nabla g(\theta) \pi(d \theta)\right)^{\top} \left(\int_{\Theta} I(\theta) \pi(d \theta)+J(\pi)\right)^{-1} \left(\int_{\Theta} \nabla g(\theta) \pi(d \theta)\right),
\end{align*}
where $J(\pi):=\int_{\Theta} \nabla_\theta \log \pi(\theta) \nabla_\theta \log \pi(\theta)^{\top} \pi(d \theta)$ is the Information theorist's Fisher information.
\end{thm}

Similar to the proof of the Cramer-Rao lower bound, we start with two lemmas.

\begin{lem}Suppose that the distribution of $X$ under $\theta$ has density $p_\theta$ with respect to some base measure $\mu$, then
\begin{enumerate}
    \item $\iint \nabla_{\theta} (p_{\theta}(x) \pi(\theta))d\mu(x) d\theta = 0$
    \item $\iint g(\theta) \cdot \nabla_{\theta} (p_{\theta}(x) \pi(\theta))d\mu(x) d\theta = - \iint \nabla_{\theta} g(\theta) \cdot p_{\theta}(x) \pi(\theta)d\mu(x) d\theta$
\end{enumerate}


\end{lem}
\begin{proof}
(a) Using $\int p_\theta(x)\,d\mu(x)=1$ and differentiation under the integral,
\[
\iint \nabla_\theta(p_\theta(x)\pi(\theta))\,d\mu(x)\,d\theta
= \int \nabla_\theta\!\left(\pi(\theta)\int p_\theta(x)\,d\mu(x)\right)\,d\theta
= \int \nabla_\theta \pi(\theta)\,d\theta.
\]
By compact support and vanishing boundary terms, $\int \nabla_\theta\pi(\theta)\,d\theta=0$.

(b) Observe that
\[
\nabla_\theta\!\big(g(\theta)p_\theta(x)\pi(\theta)\big)
= \nabla_\theta g(\theta)\,p_\theta(x)\pi(\theta) + g(\theta)\,\nabla_\theta(p_\theta(x)\pi(\theta)).
\]
Integrating over $(x,\theta)$ yields
\[
\iint g(\theta)\,\nabla_\theta(p_\theta(x)\pi(\theta))\,d\mu\,d\theta
+ \iint \nabla_\theta g(\theta)\,p_\theta(x)\pi(\theta)\,d\mu\,d\theta
= \iint \nabla_\theta\!\big(g(\theta)p_\theta(x)\pi(\theta)\big)\,d\mu\,d\theta.
\]
As in (a), differentiation under the integral gives
\[
\iint \nabla_\theta\!\big(g(\theta)p_\theta(x)\pi(\theta)\big)\,d\mu\,d\theta
= \int \nabla_\theta\!\left(\pi(\theta)g(\theta)\int p_\theta(x)\,d\mu(x)\right)\,d\theta
= \int \nabla_\theta(\pi(\theta)g(\theta))\,d\theta = 0,
\]
again by compact support/vanishing boundary terms. Rearranging yields (b).
\end{proof}

\begin{proof}[Proof of Theorem]
Define the joint density of $(X,\theta)$ (w.r.t.\ $\mu\otimes d\theta$) by
\[
m(x,\theta) := p_\theta(x)\pi(\theta),
\]
and write $\E_\pi[\cdot]$ for expectation under this joint law:
\[
\E_\pi[h(X,\theta)] := \iint h(x,\theta)\,p_\theta(x)\pi(\theta)\,d\mu(x)\,d\theta.
\]
Let
\[
U := \delta(X)-g(\theta),\qquad
s_\theta(x):=\nabla_\theta\log p_\theta(x),\qquad
\rho(\theta):=\nabla_\theta\log\pi(\theta),
\]
and define the \emph{Bayesian score}
\[
V := s_\theta(X)+\rho(\theta)=\nabla_\theta\log\big(p_\theta(X)\pi(\theta)\big).
\]
Note that $V(x,\theta)m(x,\theta) = \nabla_\theta m(x,\theta)$.

\textbf{Step 1: Identify $\E_\pi[UV]$.}
Using $V m=\nabla_\theta m$,
\[
\E_\pi[U V]
= \iint (\delta(x)-g(\theta))\,\nabla_\theta\!\big(p_\theta(x)\pi(\theta)\big)\,d\mu(x)\,d\theta.
\]
Split into two terms:
\[
\E_\pi[UV]
= \iint \delta(x)\,\nabla_\theta(p_\theta(x)\pi(\theta))\,d\mu\,d\theta
- \iint g(\theta)\,\nabla_\theta(p_\theta(x)\pi(\theta))\,d\mu\,d\theta.
\]
Since $\delta(x)$ does not depend on $\theta$, part (a) of lemma gives the first term equals $0$ and part (b) gives
\[
\iint g(\theta)\,\nabla_\theta(p_\theta(x)\pi(\theta))\,d\mu\,d\theta
= - \iint \nabla_\theta g(\theta)\,p_\theta(x)\pi(\theta)\,d\mu\,d\theta
= -\int_\Theta \nabla g(\theta)\,\pi(\theta)\,d\theta.
\]
Hence
\begin{equation}\label{eq:UV}
\E_\pi[UV] = \int_\Theta \nabla g(\theta)\,\pi(\theta)\,d\theta.
\end{equation}

\textbf{Step 2: Identify $\E_\pi[VV^\top]$.} We expand
\[
\E_\pi[VV^\top]
=\E_\pi[s_\theta(X)s_\theta(X)^\top] + \E_\pi[\rho(\theta)\rho(\theta)^\top]
+ \E_\pi[s_\theta(X)\rho(\theta)^\top] + \E_\pi[\rho(\theta)s_\theta(X)^\top].
\]
First,
\[
\E_\pi[s_\theta(X)s_\theta(X)^\top]
= \int_\Theta \E_\theta[s_\theta(X)s_\theta(X)^\top]\ \pi(\theta)\,d\theta
= \int_\Theta I(\theta)\,\pi(\theta)\,d\theta.
\]
Second, by definition,
\[
\E_\pi[\rho(\theta)\rho(\theta)^\top] = J(\pi).
\]
For the cross term, condition on $\theta$:
\[
\E_\pi[s_\theta(X)\rho(\theta)^\top]
= \int_\Theta \Big(\E_\theta[s_\theta(X)]\Big)\rho(\theta)^\top\,\pi(\theta)\,d\theta.
\]
But $\E_\theta[s_\theta(X)]=0$ since
\[
\E_\theta[s_\theta(X)]
=\int \nabla_\theta\log p_\theta(x)\,p_\theta(x)\,d\mu(x)
=\int \nabla_\theta p_\theta(x)\,d\mu(x)
=\nabla_\theta \int p_\theta(x)\,d\mu(x)
=\nabla_\theta 1 = 0,
\]
so $\E_\pi[s_\theta(X)\rho(\theta)^\top]=0$ and similarly $\E_\pi[\rho(\theta)s_\theta(X)^\top]=0$. Therefore
\begin{equation}\label{eq:VV}
\E_\pi[VV^\top] = \int_\Theta I(\theta)\,\pi(\theta)\,d\theta + J(\pi).
\end{equation}

\textbf{Step 3: A quadratic lower bound via Cauchy--Schwarz.}
For any $a\in\mathbb{R}^d$,
\[
0 \le \E_\pi\big[(U-a^\top V)^2\big]
= \E_\pi[U^2] - 2a^\top \E_\pi[UV] + a^\top \E_\pi[VV^\top]a.
\]
Minimizing the right-hand side over $a$ gives the optimizer
\[
a^\star = \big(\E_\pi[VV^\top]\big)^{-1}\E_\pi[UV],
\]
and the minimal value equals
\[
\E_\pi[U^2] - \E_\pi[UV]^\top\big(\E_\pi[VV^\top]\big)^{-1}\E_\pi[UV]\ \ge\ 0.
\]
Hence
\begin{equation}\label{eq:mainbound}
\E_\pi[U^2] \ge \E_\pi[UV]^\top\big(\E_\pi[VV^\top]\big)^{-1}\E_\pi[UV].
\end{equation}
Finally, note that $\E_\pi[U^2]=r_\pi(\delta)$ by definition, and substitute
\eqref{eq:UV}--\eqref{eq:VV} into \eqref{eq:mainbound} to obtain
\[
r_\pi(\delta)\ \ge\
\Big(\int_\Theta \nabla g(\theta)\,\pi(\theta)\,d\theta\Big)^{\!\top}
\Big(\int_\Theta I(\theta)\,\pi(\theta)\,d\theta + J(\pi)\Big)^{-1}
\Big(\int_\Theta \nabla g(\theta)\,\pi(\theta)\,d\theta\Big),
\]
which is the claimed van--Trees inequality.
\end{proof}
\begin{remark}
    Compared to the Cramer-Rao lower bound, we remove the assumption of unbiasedness by adding an extra $J(\pi)$ term.
\end{remark}


The Bayesian CRLB is a powerful tool for establishing lower bounds on the minimax risk over a local neighborhood of a parameter point $\theta_0$. 
\begin{eg}
    The \textbf{local minimax risk} is defined as:
\[
R_{\text{local-minimax}}(\theta_0, \varepsilon) = \inf_{\hat{\theta}_n} \sup_{|\theta - \theta_0| \leq \varepsilon} E_{\theta}[(\hat{\theta}_n - \theta)^2]
\]
This risk is bounded below by the Bayes risk for any prior $\pi$ supported on $[\theta_0 - \varepsilon, \theta_0 + \varepsilon]$.

Consider a prior $\nu(x)$ on $[\theta_0 - \varepsilon, \theta_0 + \varepsilon]$ by scaling a base prior $\nu_0$ on $[-1, 1]$. If $\nu_0(x) = \cos^2(\frac{\pi x}{2})$ for $|x| \leq 1$, then $J(\nu_0) = \pi^2$. Scaling this to $\nu(x) = \frac{1}{\varepsilon} \nu_0\left(\frac{x - \theta_0}{\varepsilon}\right)$ gives $J(\nu) = \frac{\pi^2}{\varepsilon^2}$. Applying the Bayesian CRLB with this prior yields:
\[
R_{\text{local-minimax}}(\theta_0, \varepsilon) \geq \left( n \int I(\theta) \nu(\theta) d\theta + \frac{\pi^2}{\varepsilon^2} \right)^{-1}
\]
\end{eg}
\begin{cor}[The Local Asymptotic Minimax (LAM) Theorem]
    By letting the local neighborhood shrink as the sample size grows (e.g., $\varepsilon = c/\sqrt{n}$), we derive the famous Hájek-Le Cam asymptotic lower bound:
\[
\lim_{c \to \infty} \lim_{n \to \infty} \inf_{\hat{\theta}_n} \sup_{|\theta - \theta_0| \leq c/\sqrt{n}} E[n(\hat{\theta}_n - \theta)^2] \geq \frac{1}{I(\theta_0)}
\]
\end{cor}

\subsection{Beyond Unbiased Estimators}
Let us first introduce a motivating example.

Consider a sample $X^{(1)}, \ldots, X^{(n)} \sim N_d(\mu, \Sigma)$ with $\Sigma=I_d$. Using the square loss $L(\mu, a)=\|\mu-a\|^2$ we easily show that the risk of any estimator $\widehat{\mu}_n$ admits the bias-variance decomposition
\begin{align*}
R\left(\mu, \widehat{\mu}_n\right)=\mathbb{E}\left\|\widehat{\mu}_n-\mathbb{E}\left[\widehat{\mu}_n\right]\right\|^2+\left\|\mathbb{E}\left[\widehat{\mu}_n\right]-\mu\right\|^2 .
\end{align*}
The MLE estimator of $\mu$ is the sample average $\bar{X}_n$. Its bias is zero, and its variance is
\begin{align*}
\mathbb{E}\left\|\frac{1}{n} \sum_i\left(X^{(i)}-\mu\right)\right\|^2=\frac{1}{n} \mathbb{E}\|\epsilon\|^2=\frac{d}{n} .
\end{align*}

In the special case when $n=1$, the MLE is $X$ which is an unbiased estimator with large variance $d$. Consider first as an alternative a simple linear estimator $\widehat{\mu}_C=C X$ with $C=\operatorname{diag}(c)$ diagonal and $c=\left(c_1, \ldots, c_d\right)$. In this case
\begin{align*}
R\left(\mu, \widehat{\mu}_C\right)=\sum_{i=1}^d\left(1-c_i\right)^2 \mu_i^2+\sum_{i=1}^d c_i^2 .
\end{align*}


Suppose we restrict ourselves to the hyperrectangular model class $\left|\mu_i\right| \leq \tau_i$. In this case we can easily find that the minimax risk is
\begin{align*}
\inf _{c \in \mathbb{R}^d} \sup _{-\tau \leq \mu \leq \tau} R\left(\mu, \widehat{\mu}_C\right)=\inf _c \sum_{i=1}^d\left(1-c_i\right)^2 \tau_i^2+\sum_{i=1}^d c_i^2=\sum_{i=1}^d \frac{\tau_i^2}{1+\tau_i^2}<d .
\end{align*}


This computation shows that for sparse model classes diagonal estimators $\widehat{\mu}_C$ may strictly dominate the MLE if $c$ is chosen carefully.

In this section, we show a famous surprising result that the observation $X=\left(X_1, \ldots, X_d\right)$ is not an admissible estimator for the parameter $\mu=\left(\mu_1, \ldots, \mu_d\right)$ unless $d \leq 2$. We start introducing the Stein's Unbiased Risk Estimates (SURE).

For a function $h: \mathbb{R}^d \rightarrow \mathbb{R}^d$ denote by $\mathrm{J} h(x)$ the Jacobian of $h$ at $x \in \mathbb{R}^d$, that is, the matrix of partial derivatives with
\begin{align*}
(J h(x))_{i j}=\frac{\partial h_i}{\partial x_j} .
\end{align*}

\begin{prop}
    [Stein's Unbiased Risk Estimates (SURE)] Let $X_1, \ldots, X_d$ be independent, $X_i \sim N\left(\mu_i, 1\right)$. Consider an estimator $\widehat{\mu}(X)$ of $\mu= \left(\mu_1, \ldots, \mu_d\right)$ and let
\begin{align*}
h(x)=x-\widehat{\mu}(x) .
\end{align*}

Suppose that $h(x)$ satisfies
\begin{enumerate}
    \item $h$ is (weakly) differentiable,
    \item $\mathbb{E}_\mu\|J h(X)\|<\infty$.
\end{enumerate}
Define
\begin{align*}
\hat{R}=d+\|h(X)\|^2-2 \operatorname{tr}(\operatorname{Jh}(X)) .
\end{align*}
Then
\begin{align*}
R(\mu, \widehat{\mu})=\mathbb{E}_\mu\|\widehat{\mu}(X)-\mu\|^2=\mathbb{E}_\mu R .
\end{align*}
\end{prop}

We start with Stein's lemma.
\begin{lem}[Stein's lemma] 
Let $X \sim N\left(\mu, \sigma^2\right)$ and let $h: \mathbb{R} \rightarrow \mathbb{R}$ be differentiable with $\mathbb{E}\left|h^{\prime}(X)\right|<\infty$. Then
\begin{align*}
E[(X-\mu) h(X)]=\sigma^2 \mathbb{E} h^{\prime}(X) .
\end{align*}
\end{lem}


\begin{proof}
   First assume $\mu=0, \sigma^2=1$. In this case we equivalently show that $\mathbb{E}(X h(X))=\mathbb{E} h^{\prime}(X)$. Without loss of generality, assume $h(0)=0$. The proof is an application of integration by parts. We have
\begin{align*}
\left(h(x) e^{-x^2 / 2}\right)^{\prime}=h^{\prime}(x) e^{-x^2 / 2}-x h(x) e^{-x^2 / 2}
\end{align*}
and so
\begin{align*}
0=\mathbb{E} h^{\prime}(X)-\mathbb{E}(X h(X)),
\end{align*}
which establishes the result. The fact that we get zero on the left follows from the fact that $\lim _{x \rightarrow \pm \infty} h(x) e^{-x^2 / 2}=0$, which can be argued since $\mathbb{E}\left|h^{\prime}(X)\right|<+\infty$. Indeed, first note that
\begin{align*}
\left|h^{\prime}(t)\right| e^{-x^2 / 2} \mathbb{1}_{[0, x]}(t) \leq\left|h^{\prime}(t)\right| e^{-t^2 / 2}
\end{align*}
and the dominating function on the right is integrable by assumption. Thus, we can use the dominating convergence theorem
\begin{align*}
\begin{aligned}
\lim _{x \rightarrow \pm \infty} h(x) e^{-x^2 / 2} & =\lim _{x \rightarrow \pm \infty} \int_0^x h^{\prime}(t) e^{-x^2 / 2} \mathrm{~d} t=\lim _{x \rightarrow \pm \infty} \int_{-\infty}^{\infty} h^{\prime}(t) e^{-x^2 / 2} \mathbb{1}_{[0, x]}(t)\d t \\
& =\int_{-\infty}^{\infty} \lim _{x \rightarrow \pm \infty} h^{\prime}(t) e^{-x^2 / 2} \mathbb{1}_{[0, x]}(t) \mathrm{d} t=0 .
\end{aligned}
\end{align*}


This establishes the result in the standard normal case. For general $\mu$ and $\sigma$ define $Z=(X-\mu) / \sigma \sim N(0,1)$. Define $\tilde{h}(z)=h(\mu+\sigma z)$. We have
\begin{align*}
\begin{aligned}
\mathbb{E}[(X-\mu) h(X)] & =\sigma \mathbb{E}(Z h(\mu+\sigma Z))=\sigma \mathbb{E}(Z \tilde{h}(Z)) \\
& =\sigma \mathbb{E} \tilde{h}^{\prime}(Z)=\sigma^2 \mathbb{E} h^{\prime}(\mu+\sigma Z)=\sigma^2 \mathbb{E} h^{\prime}(X),
\end{aligned}
\end{align*}
where moving from the first to the second line we used the case $\mu=0, \sigma=1$ proved earlier. 
\end{proof}

We need another technical lemma.
\begin{lem}
    Let $X=\left(X_1, \ldots, X_d\right)$ be a random vector with independent entries, $X_i \sim N\left(\mu_i, 1\right)$. If $h: \mathbb{R}^d \rightarrow \mathbb{R}^d$ satisfies $\mathbb{E}\|J h(X)\|<\infty$. Then
\begin{align*}
\mathbb{E}\left((X-\mu)^{\top} h(X)\right)=\mathbb{E} \operatorname{tr}(\mathrm{J} h(X))
\end{align*}
\end{lem}


\begin{proof}
    Denote by $X_{\backslash i} \in \mathbb{R}^{d-1}$ the vector $X$ with $X_i$ removed. Using the first lemma, for every $i=1, \ldots, d$
\begin{align*}
\begin{aligned}
\mathbb{E}\left(\left(X_i-\mu_i\right) h_i(X)\right) & =\mathbb{E}\left(\mathbb{E}\left[\left(X_i-\mu_i\right) h_i(X) \mid X_{\backslash i}\right]\right) \\
& =\mathbb{E}\left(\mathbb{E}\left[\left.\frac{\partial h_i(X)}{\partial x_i} \right\rvert\, X_{\backslash i}\right]\right)=\mathbb{E}\left(\frac{\partial h_i(X)}{\partial x_i}\right) .
\end{aligned}
\end{align*}
Summing over $i$ we get the result.
\end{proof}

Now the formula for the risk of Stein's Unbiased Risk Estimates follows easily.

\begin{proof}
    We have
\begin{align*}
\begin{array}{ll}
R(\mu, \widehat{\mu}) & =\mathbb{E}\left((\widehat{\mu}(X)-\mu)^{\top}(\widehat{\mu}(X)-\mu)\right. \\
& =\mathbb{E}\left((\widehat{\mu}(X)-X+X-\mu)^{\top}(\widehat{\mu}(X)-X+X-\mu)\right) \\
& =\mathbb{E}\|h(X)\|^2-2 \mathbb{E}\left((X-\mu)^{\top} h(X)\right)+\mathbb{E}\|X-\mu\|^2 \\
& =\mathbb{E}\|h(X)\|^2-2 \mathbb{E}(\operatorname{tr}(J h(X)))+\mathbb{E}\|X-\mu\|^2 \\
& =\mathbb{E} \hat{r},
\end{array}
\end{align*}
where the last equality follows because $\mathbb{E}\|X-\mu\|^2=d$, which follows because $X-\mu$ is standard normal.
\end{proof}


A natural estimator of $\mu$ is $X$ and it has constant risk $R(\mu, X)= \mathbb{E}\|X-\mu\|^2=d$. Although it is unbiased, the variance is large if $d$ is large. The James-Stein estimator of $\mu=\left(\mu_1, \ldots, \mu_d\right)$ is defined as
\begin{align*}
\delta_{J S}(X)=\left(1-\frac{d-2}{\|X\|^2}\right) X .
\end{align*}
which dominates the natural estimator when $d\ge3$.
\begin{thm}
    The risk of the James-Stein estimator is
\begin{align*}
R\left(\mu, \delta_{\mathrm{JS}}\right)=d-\mathbb{E}_\mu\left(\frac{d-2}{\|X\|}\right)^2 .
\end{align*}
In particular, the natural estimator is not admissible if $d \geq 3$.
\end{thm}
\begin{proof}
    For the James-Stein estimator, the function $h$ in SURE is
\begin{align*}
h(x)=\frac{d-2}{\|x\|^2} x
\end{align*}
and so
\begin{align*}
\hat{R}=d+\frac{(d-2)^2}{\|X\|^2}-2 \operatorname{tr}(J h(X)) .
\end{align*}
The diagonal entries of the Jacobian have a simple form
\begin{align*}
\frac{\partial h_i}{\partial x_i}=\frac{d-2}{\|x\|^2}-2 \frac{d-2}{\|x\|^4} x_i^2,
\end{align*}
giving
\begin{align*}
\operatorname{tr}(J h(X))=\left(\frac{d-2}{\|X\|}\right)^2 \quad \text { and } \quad \hat{R}=d-\left(\frac{d-2}{\|X\|}\right)^2 .
\end{align*}
Therfore, if $d \geq 3$
\begin{align*}
R\left(\mu, \delta_{\mathrm{JS}}\right)=\mathbb{E}_\mu R=d-\mathbb{E}_\theta\left(\frac{d-2}{\|X\|}\right)^2<d=R(\mu, X)
\end{align*}
proving inadmissibility.
\end{proof}


It is interesting to study the risk of a general linear estimator $\widehat{\mu}_C= C X$ for an arbitrary square matrix $C$. To use SURE, note that $h(X)= \left(I_d-C\right) X$ and so $\mathrm{J} h(x)=\left(I_d-C\right)$ giving $\operatorname{tr}(\mathrm{J} h(x))=d-\operatorname{tr}(C)$.
Therefore
\begin{align*}
R\left(\mu, \widehat{\mu}_C\right)=\mathbb{E}_\mu\left[d+\left\|\left(I_d-C\right) X\right\|^2-2 d+2 \operatorname{tr}(C)\right]=\mathbb{E}\left\|\left(C-I_d\right) X\right\|^2-d+2 \operatorname{tr}(C) .
\end{align*}


We formulate the following result without a proof.
\begin{prop}
    The linear estimator $\widehat{\mu}_C=C X$ is admissible if and only if
    \begin{enumerate}
        \item $C$ is symmetric,
        \item the eigenvalues satisfy $0 \leq \rho_i(C) \leq 1$,
        \item $\rho_i(C)=1$ for at most two $i$.
    \end{enumerate}

\end{prop}

In many situations in the expression, $C=C(\lambda)$ depends on a regularization parameter $\lambda$ so one could find optimal $\lambda^*$ by minimizing the MSE over $\lambda$. One example is given by ridge regression. 

\begin{eg}
    In this case for some $Z \in \mathbb{R}^{d \times p}$
\begin{align*}
C=Z\left(Z^{\top} Z+\lambda I_d\right)^{-1} Z^{\top}
\end{align*}
and the corresponding estimator is obtained as $\widehat{\mu}_Z=Z \widehat{\beta}$ where $\widehat{\beta} \in \mathbb{R}^p$ is obtained by solving
\begin{align*}
\min _\beta \frac{1}{2}\|X-Z \beta\|^2+\frac{\lambda}{2}\|\beta\|^2, \quad \lambda>0 .
\end{align*}
\end{eg}

\section{Lecture 4: Testing}

We start with the basic framework of hypothesis testing.

\begin{de}
   Hypothesis testing is used to determine whether the underlying parameter $\theta$ is in some specific test. Given a parameter space $\Theta$, we consider two disjoint sets $\Theta_0$ and $\Theta_1$ with the following hypothesis:
    \begin{itemize}
            \item \textbf{Null Hypothesis:} $H_0: \theta \in \Theta_0$,
            \item \textbf{Alternative Hypothesis:} $H_1: \theta \in \Theta_1$.
    \end{itemize}
   A test function is defined by a function $\phi(x)$ that maps the data to the interval $[0,1]$ that determines the outcome the test:
    \[
    \phi(x) = \begin{cases} 
    1 & \text{reject } H_0 \\
    \pi\in(0,1) & \text{reject $H_0$ with probability $\pi$} \\
    0 & \text{do not reject } H_0 
    \end{cases}
    \]
\end{de}
\begin{de}
    Let $\phi(X)$ be a test function, where $\phi(X) = 1$ means we reject the null hypothesis $(H_0)$ and $\phi(X) = 0$ means we do not reject it. We define
\begin{itemize}
    \item \textbf{Significance Level ($\alpha_{\phi}$):} The worst case of probability of a Type I error (rejecting $H_0$ when it is true).
    $$\alpha_{\phi} = \sup_{\theta \in \Theta_0} \mathbb{E}_{\theta}[\phi(X)]$$
    For a level-$\alpha$ test, we require $\alpha_{\phi} \le \alpha$ for a pre-specified $\alpha$ (e.g., $5\%$).
    \item \textbf{Power ($\beta_{\phi}(\theta)$):} The probability of correctly rejecting the null hypothesis.
    $$\beta_{\phi}(\theta) = \mathbb{E}_{\theta}[\phi(X)] \quad \text{for } \theta \in \Theta_1$$
\end{itemize}
\end{de}
\begin{remark}
    The goal is to keep $\alpha_{\phi} \le \alpha$ while maximizing $\beta_{\phi}$.
\end{remark}
\begin{de}
    A test $\phi^{\star}$ is \textbf{Uniformly Most Powerful} ($\mathbf{UMP}$) (level-$\alpha$) if for any other level-$\alpha$ test $\phi$,
$$\beta_{\phi^{\star}}(\theta) \ge \beta_{\phi}(\theta) \quad (\forall \theta \in \Theta_1).$$
\end{de}
\begin{remark}
    A UMP test is not always achievable.
\end{remark}

\subsection{Simple Testing}
\begin{de}
    For a \textbf{Simple Testing Problem}, the parameter set of the null and alternative hypotheses are singleton:
\begin{align*}
    H_0: \theta \in \Theta_0 &= \{\theta_0\} \quad (\text{under distribution } P_0) \\
    H_1: \theta \in \Theta_1 &= \{\theta_1\} \quad (\text{under distribution } P_1)
\end{align*}
\end{de}

\begin{de}
    The \textbf{Likelihood Ratio} is $L(X) = \frac{p_1(X)}{p_0(X)}$, where $p_i(x)$ is the density of the data under $P_i$.
The \textbf{Likelihood Ratio Test} is defined as:
$$
\phi(X) =
\begin{cases}
    1 & L(X) > c, \\
    \pi & L(X) = c, \\
    0 & L(X) < c.
\end{cases}
$$
\end{de}

We now state and prove the famous Neyman-Pearson lemma, which claims that to find a UMP test, it suffices to consider likelihood ratio tests.
\begin{lem}[Neyman-Pearson lemma]
    Assuming $P_0, P_1$ have common support. Given $\alpha \in (0,1)$, there exists an LRT $\phi^{\star}$ such that $\alpha_{\phi^*}=\alpha$, and any LRT $\phi^*$ such that $\alpha_{\phi^*}=\alpha$ is an UMP.
\end{lem}
\begin{proof}
   By picking $c$ and $\pi$ appropriately, we can make the likelihood ratio test $\phi^*$ a level-$\alpha$ test. Since $\phi^*$ maximizes
   $$
   \max_{\phi} \left\{ \mathbb{E}_1[\phi(X)] - c \mathbb{E}_0[\phi(X)] \right\}
   $$
   because $L(X)>c$ is equivalent to $p_1(X)>cp_0(X)$. For any other level-$\alpha$ test $\phi$, since $\phi^{\star}$ maximizes the objective, we have
\[
\mathbb{E}_1[\phi^{\star}] - c \mathbb{E}_0[\phi^{\star}]\ge \mathbb{E}_1[\phi] - c \mathbb{E}_0[\phi].
\]
Therefore,
\begin{align*}
    \mathbb{E}_1[\phi] \le \mathbb{E}_1[\phi^{\star}] + c (\underbrace{\mathbb{E}_0[\phi]}_{\le \alpha} - \underbrace{\mathbb{E}_0[\phi^{\star}]}_{= \alpha}) \le \mathbb{E}_1[\phi^{\star}]
\end{align*}
Since $\mathbb{E}_i[\phi] = \beta_{\phi}(\theta_i)$, this shows $\beta_{\phi^{\star}} \ge \beta_{\phi}$.
\end{proof}

\begin{remark}[Testing Lower Bound]
    The function that we constructed in the proof
     $$
     \mathbb{E}_1[\phi(X)] - c \mathbb{E}_0[\phi(X)] 
     $$
     is related to the lower bound of the testing error. For instance, when $c=1$, the objective function relates the minimum probability of error to the total variation distance. For any test $\phi$,
\begin{align*}
    \text{Probability of Type I + Type II error}&=\mathbb{E}_0[\phi] + \mathbb{E}_1[1 - \phi] \\& = \int (p_0(x)\phi(x) + p_1(x)(1 - \phi(x))) d\mu(x) \\
    &\ge \int \min(p_0(x), p_1(x)) d\mu(x)= 1 - d_{\text{TV}}(P_0, P_1)
\end{align*}
This minimum is achieved by the LRT with $c=1$.
\end{remark}

\begin{remark}
    Choosing $c=1$ in the Neyman-Pearson objective yields $d_{\text{TV}}$. Choosing different values of $c$ leads to an asymmetric notion of distances/indistinguishability.
\end{remark}
\begin{eg}
    Consider $X_1, \dots, X_n \sim N(\theta, 1)$ for testing $H_0: \theta=0$ vs. $H_1: \theta=\mu$. The TV distance between the $n$-sample distributions is:
$$
d_{\text{TV}}(P_1^{\otimes n}, P_0^{\otimes n}) \le \sqrt{\frac{1}{2} D_{\text{KL}}(P_1^{\otimes n} \| P_0^{\otimes n})} = \sqrt{\frac{n}{2} D_{\text{KL}}(P_1 \| P_0)}
$$
If $P_1 = N(\mu, 1)$ and $P_0 = N(0, 1)$, $D_{\text{KL}}(P_1 \| P_0) = \frac{\mu^2}{2}$.
$$
d_{\text{TV}}(P_1^{\otimes n}, P_0^{\otimes n}) \le \sqrt{\frac{n}{2} \cdot \frac{\mu^2}{2}} = \frac{\mu\sqrt{n}}{2}
$$
When $\mu = \frac{1}{\sqrt{n}}$, the testing lower bound is $1-\frac{1}{2}=\frac{1}{2}$.
\end{eg}

\begin{remark}
    To prove testing lower bounds, we often compute an upper bound on $d_{\text{TV}}$.
    The $\mathbf{d_{\text{TV}}}$ between two distributions $P_0$ and $P_1$ is defined as:
$$
d_{\text{TV}}(P_0, P_1) = \frac{1}{2} \int |p_1(x) - p_0(x)| d\mu(x)
d_{\text{TV}}(P_0, P_1) = \sup_{\|f\|_{\infty} \le 1} |\mathbb{E}_1[f(X)] - \mathbb{E}_0[f(X)]|
$$
    The fundamental Pinsker's inequality relates $d_{\text{TV}}$ and the \textbf{Kullback-Leibler Divergence} ($D_{\text{KL}}$):
    $$
    d_{\text{TV}}(P, Q) \le \sqrt{\frac{1}{2} D_{\text{KL}}(P \| Q)}.
    $$
    where the Kullback-Leibler Divergence from $Q$ to $P$ is:
$$
D_{\text{KL}}(P \| Q) = \mathbb{E}_P \left[\log \frac{dP}{dQ}\right]
$$
    The benefit of using KL divergence is that it admits the following tensorization property: for i.i.d. samples $X_1, \dots, X_n$:
$$
D_{\text{KL}}(P^{\otimes n} \| Q^{\otimes n}) = \sum_{i=1}^n D_{\text{KL}}(P \| Q) = n D_{\text{KL}}(P \| Q)
$$
\end{remark}

\begin{remark}
We can study estimation lower bounds using a testing framework. The core idea is that if two probability distributions, $P_0$ and $P_1$, are difficult to distinguish, then any estimator for a parameter $g(\theta)$ that differs under these two distributions must have a high risk.
\end{remark}

\begin{lem}[Le Cam's Lemma]
Let $\delta(X)$ be an estimator for $g(\theta)$. For any two points $\theta_0, \theta_1 \in \Theta$, the minimax risk is bounded below as follows:
$$
\inf_{\delta} \sup_{\theta \in \Theta} \mathbb{E}_{\theta}[(\delta(X) - g(\theta))^2] \ge \frac{(g(\theta_1) - g(\theta_0))^2}{8} \left(1 - d_{TV}(P_0, P_1)\right)
$$
where $P_0$ and $P_1$ are the distributions corresponding to $\theta_0$ and $\theta_1$, and $d_{TV}$ is the total variation distance.
\end{lem}

\begin{proof}
Let $\delta$ be an arbitrary test. Then, we have
\begin{align*}
    \sup_{\theta \in \Theta} \mathbb{E}_{\theta}\left[(\delta(X) - g(\theta))^2\right]&\ge \frac{1}{2} \left( \mathbb{E}_{\theta_0}\left[(\delta(X) - g(\theta_0))^2\right] + \mathbb{E}_{\theta_1}\left[(\delta(X) - g(\theta_1))^2\right] \right)
\end{align*}
Let's convert the estimation problem into a simple testing problem. Consider a test
\[
\phi(X):=1_{\left\{\left|\delta(X)-g\left(\theta_1\right)\right| \leq\left|\delta(X)-g\left(\theta_0\right)\right|\right\}}
\]
Then the test $\delta$ make an error under $\theta_i$ implies that $|\delta(X)-g(\theta_i)|\ge\frac{1}{2} |g(\theta_0)-g(\theta_1)|$. Therefore, from the testing lower bound, we have
\begin{align*}
    \frac{1}{2} \left( \mathbb{E}_{\theta_0}\left[(\delta - g_0)^2\right] + \mathbb{E}_{\theta_1}\left[(\delta - g_1)^2\right] \right)&\ge \frac{|g(\theta_0)-g(\theta_1)|^2}{8}(\text{Probability of Type I + Type II error})\\
    &\ge  \frac{|g(\theta_0)-g(\theta_1)|^2}{8}\left(1 - d_{\text{TV}}(P_0, P_1)\right).
\end{align*}
\end{proof}




\subsection{One-sided and Two-sided Testing}

We would like to generalize the conclusions in simple testing to one-sided and two-sided testing. A key property that facilitates the construction of optimal tests in these scenarios is the Monotone Likelihood Ratio (MLR).

\begin{de}[Monotone Likelihood Ratio]
A family of distributions $\{P_{\theta} : \theta \in \mathbb{R}\}$ with densities $p_{\theta}(x)$ has a \textbf{Monotone Likelihood Ratio (MLR)} in a statistic $T(X)$ if for any $\theta_1 < \theta_2$, the likelihood ratio
$$
\frac{p_{\theta_2}(x)}{p_{\theta_1}(x)}
$$
is a non-decreasing function of $T(x)$.
\end{de}

\begin{eg}
    Consider a one-parameter exponential family with density:
$$
p_{\theta}(x) = \exp(\eta(\theta)T(x) - B(\theta))h(x)
$$
The likelihood ratio is:
$$
\frac{p_{\theta_2}(x)}{p_{\theta_1}(x)} = \exp\left( (\eta(\theta_2) - \eta(\theta_1))T(x) - (B(\theta_2) - B(\theta_1)) \right)
$$
If $\eta(\theta)$ is a monotone non-decreasing function of $\theta$, then for $\theta_1 < \theta_2$, we have $\eta(\theta_1) \le \eta(\theta_2)$. The exponent is a non-decreasing function of $T(x)$, and since $\exp(\cdot)$ is increasing, the entire ratio is non-decreasing in $T(x)$. Thus, the family has MLR in $T(X)$.
\end{eg}

The MLR property is the crucial condition in the Karlin-Rubin Theorem, which guarantees the existence of a Uniformly Most Powerful (UMP) test for one-sided hypotheses.
\begin{thm}[Karlin-Rubin]
Consider the one-sided hypothesis testing problem:
$$
H_0: \theta \le \theta_0 \quad \text{vs.} \quad H_1: \theta > \theta_0
$$
If the family of distributions $\{P_{\theta}\}$ has MLR in the statistic $T(X)$, then there exists a UMP test of size $\alpha$. This test, $\phi^*(x)$, is of the form:
$$
\phi^*(x) = 
\begin{cases} 
1 & \text{if } T(x) > c \\
\gamma & \text{if } T(x) = c \\
0 & \text{if } T(x) < c 
\end{cases}
$$
where the constants $c$ and $\gamma \in [0, 1]$ are chosen to satisfy the size constraint:
$$
\sup_{\theta \le \theta_0} \mathbb{E}_{\theta}[\phi^*(X)] = \mathbb{E}_{\theta_0}[\phi^*(X)] = \alpha
$$
\end{thm}

\begin{proof}
The proof relies on the Neyman-Pearson Lemma. For any specific alternative $\theta_1 > \theta_0$, the test $\phi^*$ is the Likelihood Ratio Test (LRT) for the simple hypothesis problem $H_0': \theta = \theta_0$ vs. $H_1': \theta = \theta_1$. By the Neyman-Pearson Lemma, $\phi^*$ is the Most Powerful (MP) test for this simple problem. Since the form of the test (thresholding $T(X)$) does not depend on the specific choice of $\theta_1$ (as long as $\theta_1 > \theta_0$), the test is uniformly most powerful for all $\theta > \theta_0$.
\end{proof}

For two-sided hypothesis testing problems, such as $H_0: \theta = \theta_0$ vs. $H_1: \theta \neq \theta_0$, a Uniformly Most Powerful (UMP) test generally does not exist. A test that is optimal for alternatives $\theta > \theta_0$ is typically suboptimal for $\theta < \theta_0$. To find a meaningful "best" test in these situations, we can restrict our search to the class of \textbf{unbiased tests}.

\begin{de}[Unbiased Test]
A test $\phi$ with significance level $\alpha$ is \textbf{unbiased} if its power function, $P_{\phi}(\theta) = \mathbb{E}_{\theta}[\phi(X)]$, is never smaller than its size for any parameter in the alternative hypothesis. That is:
$$
P_{\phi}(\theta) \ge \alpha \quad \text{for all } \theta \in \Theta_1
$$
This ensures that the probability of correctly rejecting the null is always at least as large as the probability of a Type I error. For a test of $H_0: \theta = \theta_0$, this implies the power function has a minimum at $\theta_0$.
\end{de}

This concept is analogous to unbiased estimation, where we seek estimators $\hat{\theta}$ that are correct on average. The search for a Uniformly Most Powerful Unbiased (UMPU) test is the search for a test that is uniformly most powerful among all unbiased tests. Our focus for constructing such tests is typically on \textbf{exponential families}.

For a multi-parameter exponential family, the probability density function is given by:
$$
p_{\eta}(x) = \exp(\eta \cdot T(x) - A(\eta)) \cdot h(x)
$$
where $\eta$ is the natural parameter vector and $T(X)$ is the vector of sufficient statistics.
\begin{eg}
    Consider the test $H_0: \eta = \eta_0$ vs. $H_1: \eta \neq \eta_0$. For a test to be unbiased, its power function $P_{\phi}(\eta) = \mathbb{E}_{\eta}[\phi(X)]$ must have a local minimum at $\eta = \eta_0$. Assuming we can differentiate under the integral sign, this implies:
$$
\left. \frac{d}{d\eta} \mathbb{E}_{\eta}[\phi(X)] \right|_{\eta=\eta_0} = 0
$$
For an exponential family, this derivative condition is equivalent to:
$$
\text{Cov}_{\eta_0}(T(X), \phi(X)) = \mathbb{E}_{\eta_0}\left[ (T(X) - \mathbb{E}_{\eta_0}[T(X)])\phi(X) \right] = 0
$$
To construct a two-sided test, we intuitively need two thresholds, $c_1$ and $c_2$. The proposed form of the test is:
$$
\phi^*(X) = 
\begin{cases} 
1 & \text{if } T(X) < c_1 \text{ or } T(X) > c_2 \\
\gamma_i & \text{if } T(X) = c_i \text{ for } i=1,2 \\
0 & \text{if } T(X) \in (c_1, c_2)
\end{cases}
$$
The constants $c_1, c_2, \gamma_1, \gamma_2$ are found by solving the following two equations simultaneously:
\begin{enumerate}
    \item \textbf{Size constraint:} $\mathbb{E}_{\eta_0}[\phi^*(X)] = \alpha$
    \item \textbf{Unbiasedness constraint:} $\mathbb{E}_{\eta_0}[T(X)(\phi^*(X) - \alpha)] = 0$.
\end{enumerate}
This motivates the following result, whose proof is omitted here.
\end{eg}


\begin{thm}
Suppose $\eta_0$ is in the interior of the natural parameter space $\mathcal{H}$. For any $\alpha \in (0,1)$, there exists a test $\phi^*$ that satisfies the two equations above. Furthermore, this test $\phi^*$ is the UMPU test for $H_0: \eta = \eta_0$ vs. $H_1: \eta \neq \eta_0$. 
\end{thm}

\subsection{Minimax Testing}
When the parameter $\theta$ is a vector of dimension $d \ge 2$, UMP or UMPU tests often do not exist. Consider testing $H_0: \theta = \theta_0$ vs. $H_1: \theta \neq \theta_0$ for $X \sim N(\theta, \sigma^2 I_d)$.

For $d=2$, we could define different rejection regions. For example:
\begin{itemize}
    \item $\phi_1$: Reject $H_0$ when $\|X - \theta_0\|_2 > c_1$ (a circular rejection region).
    \item $\phi_2$: Reject $H_0$ when $\|X - \theta_0\|_{\infty} > c_2$ (a square rejection region).
\end{itemize}
Neither test is uniformly more powerful than the other. Test $\phi_1$ will be more powerful for alternatives along the diagonals, while $\phi_2$ will be more powerful for alternatives along the axes. We cannot compare them uniformly.


Since a uniformly optimal test may not exist, we can change the criterion for optimality. The minimax framework seeks a test that minimizes the worst-case risk. First, we redefine the hypotheses to ensure the null and alternative are well-separated.

\begin{de}[Minimax Hypotheses]
We test $H_0: \theta \in \Theta_0$ against a composite alternative that is separated by a distance $\epsilon$:
$$
H_1: \theta \in \Theta_1(\epsilon) = \{ \theta \in \Theta_1 : \text{dist}(\theta, \Theta_0) \ge \epsilon \}
$$
The distance $\epsilon > 0$ is a user's choice, representing the minimum signal strength of interest.
\end{de}

The risk of a test $\phi$ is defined as the sum of the maximum Type I and Type II error probabilities:
$$
R_{\epsilon}(\phi) = \sup_{\theta \in \Theta_0} \mathbb{E}_{\theta}[\phi(X)] + \sup_{\theta \in \Theta_1(\epsilon)} \mathbb{E}_{\theta}[1-\phi(X)]
$$
The objective in minimax testing is often to find the smallest possible separation $\epsilon$ that allows for a test $\phi^*$ to achieve a specified low risk level $\delta$. The intuition is to find the minimum signal strength required for reliable detection.

\begin{eg}[Univariate Mean] Let $X_1, \dots, X_n \stackrel{iid}{\sim} N(\theta, 1)$. We test $H_0: \theta=0$ vs. $H_1: \theta \neq 0$. The minimum separation distance required to achieve a fixed risk level $\delta$ is:
$$
\epsilon_n^* = \frac{c}{\sqrt{n}}
$$
where $c$ is a constant that depends on the desired risk $\delta$.
\end{eg} 


We consider testing a hypothesis about the mean of a multivariate normal distribution. Let $X_1, \dots, X_n \stackrel{iid}{\sim} N(\theta, I_d)$, where $\theta \in \mathbb{R}^d$. The hypothesis testing problem is:
$$
H_0: \theta = 0 \quad \text{vs.} \quad H_1(\epsilon): \|\theta\|_2 \ge \epsilon
$$
This is a problem of detecting a signal of magnitude at least $\epsilon$ against the null hypothesis of no signal.

A natural first guess for the detection boundary, $\epsilon_n$, might be based on the standard estimation error for $\theta$. The sample mean $\hat{\theta}_n = \frac{1}{n}\sum X_i$ has an expected squared error of $\mathbb{E}[\|\hat{\theta}_n - \theta\|_2^2] = d/n$. This suggests that we can only distinguish $\theta$ from $0$ when $\|\theta\|_2$ is larger than the "noise" level, perhaps leading to a conjecture that the minimax separation is $\epsilon_n \approx \sqrt{d/n}$. However, the actual rate is surprisingly faster.

\begin{thm}
The minimax testing radius for the multivariate Gaussian mean problem is
$$
\epsilon_n^* \asymp \frac{d^{1/4}}{\sqrt{n}}
$$
This means that for a constant $c_1$, if $\epsilon \le c_1 \frac{d^{1/4}}{\sqrt{n}}$, no test can reliably distinguish $H_0$ from $H_1(\epsilon)$ (the sum of errors is bounded below). Conversely, for a constant $C_1$, if $\epsilon \ge C_1 \frac{d^{1/4}}{\sqrt{n}}$, there exists a test that can reliably distinguish them (the sum of errors is bounded above).
\end{thm}
\begin{proof}
The proof has two parts.

\medskip

\noindent\textbf{Part I: Upper bound / achievability.}

Let \(Y=\frac1n\sum_{i=1}^n X_i\), then
\(
Y\sim N\left(\theta,\frac1n I_d\right).
\)
We use the statistic
\(
T=\|Y\|_2^2.
\)
Since the coordinates of $Y$ are independent, we have
\begin{align*}
\mathbb E_\theta[T]
&=
\sum_{j=1}^d \mathbb E_\theta[Y_j^2]
=
\sum_{j=1}^d\left(\frac1n+\theta_j^2\right)
=
\frac dn+\|\theta\|_2^2.\\
\operatorname{Var}_\theta(T)
&=
\sum_{j=1}^d \operatorname{Var}_\theta(Y_j^2)
=
\sum_{j=1}^d\left(\frac{2}{n^2}+\frac{4\theta_j^2}{n}\right)
=
\frac{2d}{n^2}+\frac4n\|\theta\|_2^2.
\end{align*}

Fix $a\ge 1$ and define the test
\[
\phi_a(X)
=
\mathbf 1\left\{
\|Y\|_2^2
\ge
\frac dn+a\frac{\sqrt{2d}}{n}
\right\}.
\]
The threshold is the null mean plus $a$ null standard deviations.

Under $H_0$, we have
\[
\mathbb E_0[T]=\frac dn,
\qquad
\operatorname{Var}_0(T)=\frac{2d}{n^2}.
\]
Hence, by Chebyshev's inequality,
\[
\mathbb P_0(\phi_a=1)
=
\mathbb P_0\left(
T-\mathbb E_0[T]
\ge
a\sqrt{\operatorname{Var}_0(T)}
\right)
\le
\frac1{a^2}.
\]

We now bound the Type II error. Write
\(
s=\|\theta\|_2^2,
\qquad
r=\frac{\sqrt d}{n}
\). Then the threshold is
\[
\frac dn+a\frac{\sqrt{2d}}n
=
\frac dn+a\sqrt{2}\,r.
\]
Therefore,
\[
\mathbb P_\theta(\phi_a=0)
=
\mathbb P_\theta\left(
T
<
\frac dn+a\sqrt2\,r
\right)
=
\mathbb P_\theta\left(
T-\mathbb E_\theta[T]
<
a\sqrt2\,r-s
\right).
\]
If $s>a\sqrt2\,r$, then Chebyshev's inequality gives
\[
\mathbb P_\theta(\phi_a=0)
\le
\frac{\operatorname{Var}_\theta(T)}{(s-a\sqrt2\,r)^2}
=
\frac{\frac{2d}{n^2}+\frac4n s}{(s-a\sqrt2\,r)^2}.
\]
Since $r=\sqrt d/n$, this may be written as
\[
\mathbb P_\theta(\phi_a=0)
\le
\frac{2r^2+\frac4n s}{(s-a\sqrt2\,r)^2}.
\]

Now assume
\[
s=\|\theta\|_2^2 \ge K_a r,
\qquad
K_a:=8a^2.
\]
For $s>a\sqrt2\,r$, the function
\[
s\mapsto
\frac{2r^2+\frac4n s}{(s-a\sqrt2\,r)^2}
\]
is decreasing. Therefore,
\[
\mathbb P_\theta(\phi_a=0)
\le
\frac{2r^2+\frac4n K_a r}{(K_a r-a\sqrt2\,r)^2}.
\]
Using $d\ge 1$ and $K_a=8a^2$, we obtain
\[
\frac{2r^2+\frac4n K_a r}{(K_a r-a\sqrt2\,r)^2}
=
\frac{2+\frac{4K_a}{\sqrt d}}{(K_a-a\sqrt2)^2}
\le
\frac{2+4K_a}{(K_a-a\sqrt2)^2}\le
\frac1{a^2}.
\]
Thus, we have shown that
\[
\|\theta\|_2^2
\ge
8a^2\frac{\sqrt d}{n} \implies 
\mathbb P_\theta(\phi_a=0)\le \frac1{a^2}.
\]

Combining the two bounds, we have
\[
\epsilon
\ge
2\sqrt2\,a\,\frac{d^{1/4}}{\sqrt n} \implies R_\epsilon(\phi_a)
\le
\frac1{a^2}+\frac1{a^2}
=
\frac2{a^2},
\]
This proves the upper bound.

\medskip

\noindent\textbf{Part II: Lower bound / impossibility.}

We now show that no test can succeed when $\epsilon$ is smaller than a constant multiple of $d^{1/4}/\sqrt n$.

Let \(
\delta=\frac{\epsilon}{\sqrt d}.
\)
Consider the prior $\pi$ that is uniform on the hypercube
\[
\Theta_\epsilon
=
\left\{
\theta^\omega=\delta \omega:
\omega=(\omega_1,\dots,\omega_d)\in\{-1,+1\}^d
\right\}.
\]
Every point in this prior satisfies
\[
\|\theta^\omega\|_2^2
=
\sum_{j=1}^d \delta^2
=
d\delta^2
=
\epsilon^2,
\]
so the prior is supported on the sphere $\{\|\theta\|_2=\epsilon\}$.

Let $P_\theta$ denote the joint law of $X_1,\dots,X_n$ under parameter $\theta$, and let
\[
P_\pi=\int P_\theta\,d\pi(\theta)
\]
be the corresponding mixture distribution.

For any test $\phi$, by the definition of total variation distance,
\begin{align*}
\mathbb P_0(\phi=1)
+
\sup_{\|\theta\|_2=\epsilon}\mathbb P_\theta(\phi=0)
&\ge
\mathbb P_0(\phi=1)
+
\mathbb P_\pi(\phi=0)\\
&=
1-\left(\mathbb P_\pi(\phi=1)-\mathbb P_0(\phi=1)\right)\ge
1-\|P_\pi-P_0\|_{\mathrm{TV}}.
\end{align*}
Taking the infimum over all tests,
\[
\inf_\phi
\left\{
\mathbb P_0(\phi=1)
+
\sup_{\|\theta\|_2=\epsilon}\mathbb P_\theta(\phi=0)
\right\}
\ge
1-\|P_\pi-P_0\|_{\mathrm{TV}}.
\]
Thus, it remains to show that $P_\pi$ and $P_0$ are close in total variation.

We bound total variation by the chi-square divergence:
\[
\|P_\pi-P_0\|_{\mathrm{TV}}
=
\frac12\mathbb E_0\left|\frac{dP_\pi}{dP_0}-1\right|
\le
\frac12
\left(
\mathbb E_0\left[
\left(\frac{dP_\pi}{dP_0}-1\right)^2
\right]
\right)^{1/2}.
\]
Hence
\[
\|P_\pi-P_0\|_{\mathrm{TV}}
\le
\frac12\sqrt{\chi^2(P_\pi\|P_0)}.
\]

We now compute $\chi^2(P_\pi\|P_0)$. For a fixed $\theta$, the likelihood ratio between $P_\theta$ and $P_0$ is
\[
L_\theta(X)
=
\frac{dP_\theta}{dP_0}(X)
=
\exp\left(
\theta^\top\sum_{i=1}^n X_i
-
\frac n2\|\theta\|_2^2
\right).
\]
Equivalently, since $Y=n^{-1}\sum_{i=1}^n X_i$ and $\theta^\omega=\delta\omega$,
\[
L_\theta(X)
=
\exp\left(
n\theta^\top Y
-
\frac n2\|\theta\|_2^2
\right)=
\exp\left(
n\delta\,\omega^\top Y
-
\frac n2\epsilon^2
\right).
\]
The likelihood ratio of the mixture is
\[
L_\pi(X)
=
\frac{dP_\pi}{dP_0}(X)
=
2^{-d}\sum_{\omega\in\{-1,+1\}^d} L_\omega(X).
\]
Therefore,
\[
\mathbb E_0[L_\pi^2]
=
2^{-2d}
\sum_{\omega,\omega'}
\mathbb E_0[L_\omega L_{\omega'}].
\]
Under $P_0$, $Y\sim N(0,n^{-1}I_d)$. Hence, for fixed $\omega,\omega'$,
\[
\mathbb E_0[L_\omega L_{\omega'}]
=
\exp\left(
n\delta^2 \omega^\top\omega'
\right).
\]
Indeed, this follows from the Gaussian moment generating function.

Thus,
\[
\mathbb E_0[L_\pi^2]
=
2^{-2d}
\sum_{\omega,\omega'}
\exp\left(
n\delta^2 \sum_{j=1}^d \omega_j\omega_j'
\right).
\]
For independent uniformly distributed signs $\omega,\omega'$, each product
$\omega_j\omega_j'$ is again a uniform sign. Therefore,
\[
\mathbb E_0[L_\pi^2]
=
\left(
\frac{e^{n\delta^2}+e^{-n\delta^2}}2
\right)^d
\]
Since $\delta^2=\epsilon^2/d$, using the elemenetary inequality $(e^x + e^{-x})/2 \le e^{x^2/2}$, we obtain
\[
\chi^2(P_\pi\|P_0)
=
\mathbb E_0[L_\pi^2]-1
\le
\exp\left(
\frac d2\left(\frac{n\epsilon^2}{d}\right)^2
\right)-1
=
\exp\left(
\frac{n^2\epsilon^4}{2d}
\right)-1..
\]
Choose $c_1>0$ small enough, we have
\[
\epsilon
\le
c_1\frac{d^{1/4}}{\sqrt n} \implies \frac{n^2\epsilon^4}{2d}
\le
\frac{c_1^4}{2}\implies \chi^2(P_\pi\|P_0)\le \frac49
\]
Therefore,
\[
\|P_\pi-P_0\|_{\mathrm{TV}}
\le
\frac12\sqrt{\frac49}
=
\frac13.
\]
and hence,
\[
\inf_\phi
\left\{
\mathbb P_0(\phi=1)
+
\sup_{\|\theta\|_2=\epsilon}\mathbb P_\theta(\phi=0)
\right\}
\ge
1-\frac13
=
\frac23.
\]
Since the sphere $\{\|\theta\|_2=\epsilon\}$ is contained in the alternative
$\{\|\theta\|_2\ge \epsilon\}$, the same lower bound also holds for the composite alternative
$H_1(\epsilon):\|\theta\|_2\ge \epsilon$.

Combining the upper and lower bounds, we obtain
\[
\epsilon_n^*\asymp \frac{d^{1/4}}{\sqrt n}.
\]
\end{proof}




